% =====================================================================
%  Capitolo 3 — Il ciclo di produzione, l'iterazione e la documentazione
%  Fonte: trascrizione della lezione GD3 (Game Design, primo anno)
%  Pacchetti richiesti nel preambolo del libro:
%    \usepackage[italian]{babel}, \usepackage[T1]{fontenc},
%    \usepackage[utf8]{inputenc} (solo pdfLaTeX), \usepackage{booktabs}
% =====================================================================

\chapter[Iterare e documentare]{Iterare e documentare: il ciclo di produzione e le regole della documentazione}
\label{cap:produzione-documentazione}

I primi due capitoli hanno descritto il mestiere del designer e il modo in
cui un designer guarda a un gioco. Questo terzo capitolo apre un modulo
nuovo, quello dedicato alla \textbf{documentazione e al pitching}: come si
documentano le cose, come le si spiega, come le si comunica. Si parla, ancora
una volta, di \emph{linguaggio e comunicazione}: il direttore aveva definito
il documento come la forma di comunicazione ``pratica'', condensata, del
designer, ed è di questo che ci si occupa.

La lezione è tenuta da Mattia Podini, docente di game design della sede di
Milano, che si definisce un appassionato di documentazione (``è una mia
specialità''). Il percorso del capitolo ne ricalca la struttura:
\begin{enumerate}
  \item una sessione di domande e risposte sulla lezione precedente, che
        chiarisce e amplia alcuni punti dei framework di Schell, di Romero e
        Schreiber e di Bura (sezione~\ref{sec:qa-gd2});
  \item un giro nel Drive del corso e le regole di consegna dei task
        (sezione~\ref{sec:drive-regole});
  \item la lezione ``zero'' di documentazione: il \textbf{ciclo di
        produzione} e il principio dell'\textbf{iterazione}, con i suoi
        corollari sul target audience, sui test e sui pilastri del progetto
        (sezioni~\ref{sec:ciclo-produzione}, \ref{sec:domande-iterazione}
        e~\ref{sec:principi-iterazione});
  \item la prima vera lezione di documentazione: perché documentare, le
        \textbf{golden rules} della buona documentazione, come scegliere lo
        strumento giusto e le regole formali del corso
        (sezioni~\ref{sec:perche-documentare}, \ref{sec:golden-rules},
        \ref{sec:scegliere-documento} e~\ref{sec:regole-formali});
  \item il task della settimana (sezione~\ref{sec:task-gd3}).
\end{enumerate}

Già in apertura il docente avverte che la lezione sarà un ``grande spasso''
e precisa, a scanso di equivoci, che anche \emph{spasso} è una parola
proibita (si veda la sezione~\ref{sec:vocabolario}).

% ---------------------------------------------------------------------
\section{Domande e risposte sulla lezione precedente}
\label{sec:qa-gd2}

Prima di entrare nel vivo, la lezione si apre con un ripasso guidato dalle
domande degli studenti. Il docente avverte che le slide della lezione
precedente sono state \textbf{aggiornate}: ha aggiunto un paio di cose e ne ha
rimodernate altre, quindi conviene riguardarle.

\subsection{Schell: game piece, verbi e regole}

La prima domanda (Valerio) riguarda il blocco delle \emph{meccaniche} nel
framework di Jesse Schell (sezione~\ref{sec:schell}): qual è la differenza
tra \emph{game piece} e \emph{verbi di gioco}?

\begin{description}
  \item[Game piece (pezzo di gioco)] ha dentro di sé delle statistiche,
        delle informazioni pubbliche e private, una serie di caratteristiche.
        Va considerato come un \textbf{oggetto di gioco che cambia stato o fa
        cambiare stato ad altre cose}. In \emph{Super Mario} un pezzo di gioco
        può essere un nemico (un Goomba), ma anche il fiore di fuoco, la
        moneta, o Mario stesso.
  \item[Verbo di gioco] ha generalmente a che fare \textbf{solo con il
        giocatore}: serve al designer per definire che cosa il giocatore può
        fare (ed eventualmente che cosa non può fare). I verbi di Mario sono
        saltare, sparare (se ha il fiore di fuoco), camminare, corricchiare e
        correre: premendo leggermente la levetta si cammina, poi si inizia a
        corricchiare e infine si corre, una gradazione tipica dei giochi
        action (lo fa anche \emph{Uncharted}). Il verbo deve essere qualcosa
        di \textbf{attuabile dal giocatore} che si traduce in un verbo, in
        un'azione.
  \item[Regole] sono \textbf{tutto ciò che non ricade negli altri due
        blocchi}. Il limite dello spazio; il fatto che in \emph{Mario}, se
        finisce il timer, è game over; la condizione per finire un livello
        (raggiungere la bandiera, il castello). Tutto ciò che fa da
        \emph{struttura} del sistema.
\end{description}

\paragraph{Il verbo di gioco fuori dal framework.} Al di là del framework, il
termine ``verbo di gioco'' assume un \textbf{senso più generale} quando si
parla in contesti meno strutturati: in un brainstorming con un altro
designer, mentre si decide qualcosa, o quando si deve far capire il proprio
gioco, la propria visione, a qualcun altro. In \emph{Call of Duty}, come in
qualsiasi sparatutto, il verbo principe è \emph{sparare}. Ma poi, da gioco a
gioco, bisogna capire che cosa significa: il verbo \emph{sparare} funziona sia
per \emph{Call of Duty} sia per \emph{Portal}, solo che in \emph{Call of Duty}
si spara per uccidere i nemici, in \emph{Portal} si spara (con la Portal Gun)
per risolvere puzzle. Nel mondo moderno il termine si usa quindi in maniera
molto generica; quanto generica, lo si capirà nel tempo.

\subsection{Romero e Schreiber: i player sono il target audience}

La seconda domanda riguarda il framework di Brenda Romero e Ian Schreiber
(sezione~\ref{sec:romero-schreiber}): è corretto intendere i \emph{player}
come il target? Sì: \textbf{i player sono esattamente il target audience}.

Qui il docente coglie l'occasione per un punto che non era emerso abbastanza:
molti dei termini dei framework (meccaniche, estetiche, player\dots)
\textbf{hanno molteplici definizioni a seconda del framework}. Alcuni sono
linguaggi tecnici riutilizzati con accezioni diverse. Nel framework di Romero
e Schreiber i player sono il target audience e \textbf{vanno definiti}: non è
sufficiente scrivere ``casual gamer''. Servono informazioni specifiche:
\begin{itemize}
  \item nazionalità ed eventualmente etnia;
  \item sesso o genere in cui si identificano;
  \item età;
  \item potenzialmente anche il background.
\end{itemize}
In generale, \textbf{tutti i dati che si possono ricavare}, purché si tratti
di dati \emph{che hanno un riscontro}. Come si fa questa ricerca verrà
spiegato più avanti nel corso.

\subsection{Perché due framework di core elements}

Il motivo per cui sono stati presentati due framework di \emph{core elements}
non è solo dare più strumenti (che resta la ragione principale). Nel piccolo
panorama dell'industria dei videogiochi, per quanto giovane, il framework di
Schell è stato \textbf{più condiviso di recente}. C'è stato un periodo, quello
delle console dei primi anni Duemila, in cui si preferiva Romero e Schreiber.
Perché? Si immagini di giocare uno stesso titolo, ad esempio \emph{Halo}, su
due console della stessa generazione, come Xbox 360 e PlayStation 3: che
differenza ci sarebbe? Probabilmente nessuna. Il pad è generalmente lo stesso
e la tecnologia è simile, quindi era più utile definire un gioco attraverso le
\emph{game qualities}, cioè elementi atomici.\footnote{L'esempio è ipotetico:
\emph{Halo} è una serie esclusiva Microsoft, e solo di recente è stato
annunciato il suo arrivo su altre piattaforme, cosa a cui si allude in
lezione parlando di un \emph{porting}.}

Oggi invece giocare da mobile su uno schermo, o su un Nintendo DS con uno
schermo interattivo, rispetto a qualcos'altro, apre possibilità di design
diverse da quelle di altri progetti. Se la \textbf{tecnologia è un elemento
dominante}, può essere più rilevante usare il framework di Schell (che la
include esplicitamente) piuttosto che quello di Romero e Schreiber. È un primo
esempio concreto del principio delle lenti (sezione~\ref{sec:lenti}): la
scelta del framework non dipende solo dalla propria sensibilità, ma anche da
ciò che si ha davanti.

\paragraph{Quale framework per il task.} Per il task in corso
(sezione~\ref{sec:task-gd2}) il consiglio è di usare uno dei due framework di
core elements, \textbf{Schell} oppure \textbf{Romero e Schreiber}. Si può
usare anche quello di Stéphane Bura (il docente, scherzando, ammette di non
ricordarne bene il nome e lo fa sillabare in aula: B-U-R-A, S-T-E-P-H-A-N-E),
ma è più complesso. Con una nota di ironia, il docente aggiunge che il
direttore ``sbaglia parecchie nomenclature, ma se le ricorda'': il consiglio
è di seguire il filo del suo discorso anche quando i nomi non sono esatti.

\subsection{Come si usa la matrice di Bura}

La tabella vista nella lezione precedente (tabella~\ref{tab:bura}) è una
\textbf{versione semplificata} della matrice di Bura, che si può trovare
facendo ricerca. Il docente ne completa la spiegazione.

\paragraph{La premessa.} L'assunto da cui parte Bura è che gli elementi della
matrice possano essere combinati e ricombinati fino ad avere una
\textbf{formula} quasi matematica e \textbf{riproducibile}: per creare un
certo tipo di gioco, o per far evocare un certo tipo di emozioni, in maniera
riproducibile.

\paragraph{La conseguenza: dalla matrice al design.} Finora i core elements
sono serviti come \textbf{metodo di analisi} (ed è questo il task in corso);
più avanti andranno usati anche come \textbf{metodo di design}. Con Bura si
prende una meccanica, un pezzo di gioco, e si trova \textbf{la cella esatta}
della matrice in cui ricade. Che cosa succede se quella cella viene
\emph{potenziata} o \emph{bloccata}, poco o tanto? Qual è il responso
emotivo? Esempio: in un gioco horror, se si inizia a togliere
\textbf{freedom} all'\textbf{action}, ad esempio lato level design, e si
stringono i corridoi fino a lasciare al giocatore un solo corridoio lineare,
il giocatore (``perdonate il francese'') inizierà ad avere molta paura. O
almeno questo è il presupposto, ciò che si spera; e non è l'unico elemento in
gioco. Questa è la parte finale del framework, su cui eventualmente si
tornerà in futuro.

\paragraph{Il ciclo variabili--emozioni--interazioni.} Un altro punto legato
alla premessa: si può partire da dove si vuole, ma le \textbf{variabili di
gameplay} (che si traducono, o almeno si identificano, nelle celle del
framework) influenzano le emozioni del giocatore; il giocatore interagisce in
qualche modo con l'ambiente di gioco, e l'ambiente di gioco, a sua volta,
cambia le variabili di gameplay. Se si ha molta vita si fa ``lo smargiasso'',
si va tronfi e senza paura nella mischia; combattendo si perde vita, e allora
si diventa più cauti.

\paragraph{Esempio: Into the Breach.} Il docente consiglia caldamente
\emph{Into the Breach} (``dovete giocarlo''; tra l'altro è stato distribuito
gratuitamente su Epic Games ed è incluso nell'abbonamento Netflix): una sorta
di \emph{Pacific Rim} su griglia, in cui si combatte contro i Kaiju con tre
unità (i \emph{mech}), ciascuna con un pilota. Ne nasce in aula un'accesa
discussione su quanta casualità ci sia nel gioco. La posizione del docente è
che il gioco, pur essendo un roguelike (casuali sono i power-up, la creazione
delle missioni, i biomi da affrontare), \textbf{in battaglia} sia quasi del
tutto privo di casualità, ``come gli scacchi'': un PvE asimmetrico, con i
comportamenti dei nemici dichiarati in anticipo, che dipendono dalla posizione
dei mech, dalla distanza e dal massimo danno possibile. L'unica eccezione è la
\textbf{griglia energetica}.\footnote{In \emph{Into the Breach} le intenzioni
dei nemici vengono mostrate prima che agiscano, rendendo ogni turno un
problema a informazione quasi completa; la principale fonte di casualità
durante il combattimento è la \emph{grid defense}, la probabilità che un
edificio resista a un attacco, che parte dal 15\% e può essere aumentata.
Alcuni studenti fanno notare che c'è comunque della casualità nelle scelte dei
nemici; il docente rimanda ai talk degli sviluppatori.}

Il punto dell'esempio, però, è il ciclo emozioni--interazioni. Se un mech
finisce le tacche di vita si può perdere il pilota, e i piloti sono
generalmente più forti dell'IA che comanda il mech in loro assenza (hanno
delle abilità passive). La griglia dà dei bonus, ma se arriva a zero è game
over: è, di fatto, la vita del giocatore. Allora: se si ha molta griglia si
può rischiare di farla cadere (``500 civili in meno che saranno, un palazzo
demolito che sarà, per il bene dell'umanità''); se si ha una sola tacca, si
sarà molto più cauti. Oppure, e qui entrano in gioco i \emph{tipi di
giocatori}, di cui si parlerà più avanti, si può rischiare ugualmente, perché
c'è il 15\% di possibilità che l'attacco a un edificio non vada a segno: ``me
la gioco'', e magari va bene. Le variabili modificano il loop:
\begin{itemize}
  \item \textbf{per il giocatore}, il ciclo va dalle variabili alle emozioni
        alle interazioni;
  \item \textbf{per il designer}, il ciclo si percorre \emph{a ritroso}: si
        parte dalle emozioni che si vogliono far provare (la lista vista nella
        lezione precedente) e si pensa a quali variabili di gameplay e quali
        interazioni possono influire su di esse.
\end{itemize}
A volte questi elementi vengono scambiati di posto: dipende dal framework e
dal design.

\paragraph{Le scelte come una forbice.} Un ultimo punto: sia nel prendere una
scelta sia nel mappare un gioco dall'inizio, le scelte diventano
\textbf{progressivamente più ampie dall'inizio alla metà} del gioco e
\textbf{progressivamente più strette dalla metà alla fine}. In
\emph{Pokémon} all'inizio si può scegliere tra tre starter; poi nel percorso
se ne possono catturare altri cinque, poi altri centocinquantuno, e così via.
Ma \emph{Pokémon} ha un finale: si batte la Lega (i Superquattro), ``baci e
abbracci''. Non vale per tutti i giochi: nei \emph{sandbox} la forbice è
molto ampia e potenzialmente non si chiude mai.

\paragraph{Altri esempi delle celle.} Il docente mostra anche gli esempi che
aveva preparato per le celle e che non c'era stato tempo di vedere:
\begin{description}
  \item[Freedom $\times$ Action] la chiave che apre la porta.
  \item[Freedom $\times$ System] l'esplorazione: rischio di andare al
        prossimo falò, sapendo che in mezzo ci sono un certo numero di strade
        ed eventuali nemici pronti a ``massacrarmi di botte''.
  \item[Freedom $\times$ Self] non è solo quanto detto nella lezione
        precedente: è anche avere la \textbf{capacità di attuare ciò che si
        vuole in quanto individuo}, avere gli strumenti per farlo. In
        \emph{Celeste} (un platform con una colonna sonora notevole) si può
        attivare dalle opzioni un cronometro che mostra quanto tempo si sta
        impiegando a finire il gioco: uno strumento pensato per gli
        \emph{speedrunner}, che trovano all'interno del gioco ciò che serve a
        fare quello che vogliono. Funziona anche a bassissimo livello: avere
        nel gioco gli strumenti per personalizzare i propri controlli, se lo
        si ritiene una necessità.
  \item[Freedom $\times$ Social] poter parlare con altre persone.
\end{description}

\subsection{Una critica personale al framework di Bura}

Il docente premette di essere umano e quindi di avere delle preferenze: non è
un grande fan del framework di Bura, e farà emergere le criticità dei
materiali quando ci saranno.
\begin{itemize}
  \item Termini come \emph{freedom} non sono davvero adottati nel design, se
        non in maniera gergale; esistono sinonimi.
  \item \emph{Mastery} è, nella sua visione, anche un \textbf{responso
        emotivo}: lo sviluppo di competenze è una possibile emozione, e
        comparirà nella lista. Esempio: il docente è ``scarsissimo'' a
        scacchi e ha amici molto forti che lo sfidano continuamente. Se
        volessero davvero dimostrare la propria maestria, che cosa farebbero?
        Andrebbero contro persone forti: nei tornei si cerca di avere un
        \emph{bracket} bilanciato, in modo da affrontare giocatori con la
        stessa abilità o con abilità superiore, ed è lì che si vede chi è il
        migliore. Quindi il loro obiettivo, molto probabilmente, è un altro:
        farlo sentire inadeguato. Sono due risposte emotive molto diverse, e
        per questo il docente colloca la \emph{mastery} altrove.
  \item Lo stesso vale per la \emph{freedom}: ogni gioco ha dei limiti,
        \textbf{la vera libertà non esiste} (ci sarà una lezione apposita).
\end{itemize}

\paragraph{Lo stesso termine, significati diversi: le estetiche.} Vale anche
per le \emph{estetiche}. Nel framework di Schell includono la
\emph{vestizione}, il \emph{dressing}: l'arte, la direzione artistica, i
suoni, tutto ciò che viene apposto sopra le meccaniche. In un altro framework,
invece, si parla di \textbf{game aesthetics}, che sono un'altra cosa: sono
quelle che compaiono nella definizione di gioco del docente
(sezione~\ref{sec:non-definibile}), e che per ora vanno lette come
sinonimo di \emph{emozioni}.\footnote{Il riferimento implicito è con ogni
probabilità al framework MDA (\emph{Mechanics, Dynamics, Aesthetics}) di
Hunicke, LeBlanc e Zubek, in cui le \emph{aesthetics} sono le risposte
emotive desiderate nel giocatore.}

\paragraph{Parlare la stessa lingua.} Esistono quindi molti sinonimi, alcuni
più condivisi di altri. Non significa che uno sia giusto e l'altro
sbagliato: dipende dal contesto, dall'azienda in cui si lavora, da chi si
collabora. L'obiettivo è \textbf{comunicare nella stessa lingua}, e a volte,
ahimè, bisognerà piegarsi agli altri, perché un termine è condiviso o
imposto dall'alto.

\paragraph{Aneddoto: lo stagger.} Il docente racconta (senza dire dove né con
chi) un caso che ``ha fatto impazzire'' tutto il team di design. In una grossa
azienda, in un gioco con un sistema di combattimento, il team presentava la
documentazione al creative director in sessioni di \emph{briefing} e
\emph{debriefing}. Generalmente, quando un nemico è \emph{staggerato} (in
\emph{Sekiro} l'equivalente è lo sbilanciamento), non viene solo stordito: lo
si interrompe, è un termine specifico. Il creative director voleva invece che
\emph{stagger} indicasse un'altra cosa, traslando termini come \emph{stun} e
\emph{stagger} in accezioni non di uso comune, e non legate a come il
giocatore le percepisce. Il team ha dovuto rispettare l'ordine, perché veniva
dall'alto; lo ha fatto presente, ma la segnalazione non è stata accolta, e
questo ha creato non pochi problemi. È un problema di direzione; ma dimostra
che l'obiettivo è trovare il linguaggio più comune a tutti, entro dei limiti:
un'azienda che assume non nasce con il nuovo assunto, ha già una struttura e
delle regole, e bisognerà stare a quelle (limiti che, peraltro, possono anche
essere utili).

\subsection{Scrivere da designer, non da recensori}

Una domanda di Federica sulla lezione precedente (``bello'', ``brutto'',
``noioso'') porta a chiarire il senso dell'obbligo sul vocabolario
(sezione~\ref{sec:vocabolario}): serve a far \textbf{parlare} gli studenti da
designer, e in futuro a farli \textbf{ragionare} da designer. Il docente sa
già che succederà: molti inizieranno a scrivere \emph{come recensori}, non
come designer, perché si è abituati a un certo tipo di editoria. Dire che un
gioco è bello o brutto, che piace o non piace, è legato a una sensibilità
personale che non deve esserci, perché \textbf{il design è \emph{egoless}}
(sezione~\ref{sec:bateman}).

Un modo per togliere questa sensibilità sono i \textbf{dati}. È il discorso
dei player di Romero e Schreiber: colloquialmente si può dire che
\emph{Candy Crush} è un gioco ``adatto alle casalinghe di mezza età'', ma nei
documenti questa frase non esiste; al massimo nel brainstorming, se va bene.
Un documento tecnico richiede dati, e deve essere \textbf{il più inconfutabile
possibile}, per due ragioni:
\begin{itemize}
  \item il resto del team si aspetta che si sia fatto un certo tipo di
        ricerca e che si sia sicuri della propria idea;
  \item il designer stesso, per primo, deve smontare la propria idea ed
        essere sicuro dei dati: è il metodo scientifico.
\end{itemize}
Le cose poi cambiano, le meccaniche si iterano, ma l'assunto di base è
essere il più sicuri possibile. Il resto del team viene a chiedere al
designer quando ha un dubbio: se nemmeno il designer è sicuro di ciò che
scrive e dice, come si fa?

% ---------------------------------------------------------------------
\section{Il Drive del corso e le regole dei task}
\label{sec:drive-regole}

\subsection{La struttura del Drive}

Tutti gli studenti sono stati aggiunti al Drive del corso, e il docente ne fa
un giro esplorativo.
\begin{itemize}
  \item C'è la cartella del \textbf{primo anno di game design}; ci saranno poi
        quelle del secondo e del terzo.
  \item Dentro, le cartelle sono divise in \textbf{lezioni}, \textbf{materiale
        studenti} e \textbf{task}.
  \item Le lezioni sono divise in \emph{Game Design \& Method} e
        \emph{Documentation} (questa lezione appartiene alla seconda), più le
        lezioni condivise con i dipartimenti coinvolti (ad esempio quelle con
        narrativa).
  \item \textbf{Materiale studenti} è pensato per documenti aggiuntivi caricati
        dagli studenti: ad esempio, se qualcuno trovasse il PDF del framework
        di Bura, lo si potrebbe caricare lì. L'anno precedente non è servito a
        nulla, perché nessuno ha caricato niente, ma si cerca di mantenerlo.
        Materiali come i pacchetti di asset sono invece caricati di solito dal
        docente di level design nella cartella di level design.
\end{itemize}

\paragraph{I permessi.} Gli studenti hanno i permessi per \textbf{modificare
tutto}: possono rinominare, aggiungere cartelle e così via (``se volete
distruggere casa vostra, fatelo, ma non ve lo consiglio''). Non possono però
\textbf{cancellare}: se per errore si crea un documento che non dovrebbe
esserci, bisogna avvisare il docente, che lo cancellerà. Senza abusarne:
prima di creare documenti, assicurarsi che siano quelli corretti.

\subsection{Come si consegna un task}

\begin{enumerate}
  \item Per ogni task \textbf{uno solo} degli studenti (il primo che arriva si
        prende ``questo onere e onore'') crea una cartella dentro la cartella
        dei task della disciplina (Task GD, Task ND o Task LD). Per il task
        degli arcade la cartella è già stata creata dal docente.
  \item La cartella si rinomina nel formato \textbf{anno mese giorno titolo
        del task}. Chi la crea ha il merito di scegliere il titolo, purché
        abbia senso (niente ``uova di Pasqua'' se il task non riguarda quello).
  \item Dentro quella cartella ciascuno crea \textbf{direttamente} il proprio
        documento.
  \item Se si devono caricare più documenti, si crea una sottocartella
        \textbf{cognome e nome} e li si carica lì. Se il documento è uno solo,
        la sottocartella non serve.
\end{enumerate}

\paragraph{Scadenza.} Si ha \textbf{una settimana} dall'assegnazione del
task. Il task in corso è stato assegnato mercoledì, quindi va consegnato
\textbf{entro mercoledì alle ore 13}. Se lo si carica dopo, il docente assume
che non sia stato fatto; se lo si modifica dopo la scadenza, il docente non lo
legge. Solo dopo la scadenza il docente inizia a commentare.

\paragraph{Durante la revisione.} Se arrivano le notifiche di Google e si vede
che il docente è dentro il documento e sta commentando, non bisogna
``andargli sopra'', mandare emoji o fare domande: si aspetta che abbia finito,
perché altrimenti lo si rallenta (e ci si mette ulteriormente in ansia). Il
docente ammette di esserci passato anche lui: la curiosità è tanta e
l'emozione molta. Come sapere quando ha finito? La cartella del task
\textbf{sparisce da Task GD e finisce in ``Revisionati''}. Da quel momento si
ha piena libertà di leggere i propri documenti, rispondere ai commenti e fare
domande. Per i dubbi si può sempre parlare prima o fuori dalla lezione, o
scrivere su Discord. Ma prima di contestare un commento (il docente può
fraintendere qualcosa) bisogna \textbf{leggere i commenti di tutti}, non solo
i propri.

\paragraph{Solo la suite di Google.} Per ora, e per il prossimo task, si
possono usare \textbf{solo gli strumenti della suite di Google}. È un
consiglio che per ora vale come obbligo, e molte di queste regole verranno
meno nel tempo (il perché si vedrà in questa stessa lezione). Il documento va
\textbf{creato direttamente nella cartella del task}: non lo si crea nel
proprio Drive per poi copiarlo lì alla fine. Se si vuole una copia personale,
la si crea a partire dall'originale, ma l'originale deve stare nella cartella
del corso. Le ragioni:
\begin{itemize}
  \item nessuno ruba i documenti (all'inizio ci sta un po' di gelosia), e
        bisogna abituarsi al fatto che \textbf{i documenti che si producono
        non sono propri, ma dell'azienda};
  \item la cronologia delle modifiche permette al docente, in caso di dubbi,
        di verificare che il task non sia stato fatto la sera prima della
        scadenza;
  \item serve a tutela degli studenti in caso di copiature, che si assume non
        ci siano mai ma che ogni anno capita di scoprire; la punizione è
        estremamente dura.
\end{itemize}
Tutte queste regole sono raccolte in un documento apposito nella cartella
dei task, che verrà aggiornato man mano che se ne aggiungeranno di nuove (si
veda anche la sezione~\ref{sec:regole-formali}).

\paragraph{La punizione della parola proibita.} Lo studente che deve
consegnare due analisi (sezione~\ref{sec:task-gd2}) chiede se debba
analizzare due giochi con un solo framework. Il consiglio è: \textbf{un gioco
analizzato con Romero e Schreiber e un gioco diverso analizzato con Schell},
quindi due giochi e due framework diversi, sconsigliando Bura per semplicità
(chi vuole impegnarsi di più può comunque farlo). In questo caso la punizione
coincide con il task del giorno, ma generalmente le punizioni sono altro,
qualcosa di diverso dall'argomento in corso.

% ---------------------------------------------------------------------
\section{Il ciclo di produzione}
\label{sec:ciclo-produzione}

\subsection{La lezione zero}

La prima lezione del modulo di documentazione è la lezione \textbf{zero}: non
la più importante, ma una di quelle cose che devono entrare nella
\emph{forma mentis}. Si parla di \textbf{ciclo di produzione}. (Come da
tradizione, le lezioni di documentazione si aprono sempre con un
\emph{meme}, che il docente aggiorna ogni anno; quest'anno uno su
\emph{Metroid Prime 4}.)

Il primo ciclo mostrato è, per ammissione del docente, \textbf{bruttissimo}:
fatto con Paint. E c'è una ragione, che si capirà più avanti
(sezione~\ref{sec:segreti-iterazione}).

\subsection{Preproduzione, produzione, post-produzione}

Dal secondo anno (non al primo, per le ragioni spiegate nella prima lezione)
si vedranno diverse metodologie di produzione, tra cui quella principale. In
generale, però, un progetto si divide in tre grandi blocchi.

\begin{description}
  \item[Preproduzione] il blocco iniziale, in cui si fa \textbf{tanta ricerca,
        tanta documentazione}, tante cose in vista del momento in cui ``ci si
        sporcherà le mani''. Le tre fasi non occupano un terzo del tempo
        ciascuna: questa è probabilmente la più breve, purtroppo. Purtroppo
        perché \textbf{quanto più è lunga e profonda la preproduzione, tanto
        meglio usciranno le altre due}. È una cosa che la gente non vuole mai
        capire, nemmeno le aziende; il docente capisce che dal punto di vista
        dei costi sia la fase meno conveniente, perché non sempre ha un
        \emph{output} vendibile, ma assicura che è una delle più importanti,
        soprattutto per partire con il piede giusto.
  \item[Produzione] la fase che prende gran parte del tempo: quando si fanno
        le cose, quando il gioco si crea.
  \item[Post-produzione] a volte una fetta molto breve: tutto ciò che si fa
        dopo il rilascio o in vista del rilascio, cioè \emph{polish} estremo,
        \emph{bug fixing} ed eventualmente la creazione di DLC.
\end{description}

Il ciclo include anche il \textbf{testing}, la \textbf{beta} (un tipo di
versione: ce ne sono diverse, se ne vedranno altre) e il \textbf{rilascio}
(\emph{release}).

\subsection{DevOps e LiveOps}

Questa suddivisione viene da altre industrie: preproduzione, produzione e
soprattutto post-produzione esistono anche nell'industria cinematografica. Per
i videogiochi, però, da qui in avanti (``nei secoli dei secoli, amen'') il
termine \emph{post-produzione} si può dimenticare, perché esistono termini
più corretti:

\begin{table}[htbp]
  \centering
  \begin{tabular}{@{}p{0.18\textwidth}p{0.36\textwidth}p{0.36\textwidth}@{}}
    \toprule
    & \textbf{DevOps} & \textbf{LiveOps} \\
    \midrule
    \textbf{Quando} & tutto ciò che avviene \emph{prima} del rilascio
                    & tutto ciò che avviene \emph{dopo} il rilascio \\
    \addlinespace
    \textbf{Tipico di} & giochi che escono ``finiti''
                       & \emph{Games as a Service}: giochi mantenuti in vita
                         dopo il rilascio \\
    \addlinespace
    \textbf{Esempi} & \multicolumn{1}{c}{---}
                    & \emph{Candy Crush}, \emph{Clash Royale},
                      \emph{Fortnite}, \emph{Apex Legends}, gran parte dei
                      giochi mobile \\
    \bottomrule
  \end{tabular}
  \caption{DevOps e LiveOps nel ciclo di vita di un gioco.}
  \label{tab:devops-liveops}
\end{table}

\paragraph{Esempio: Apex Legends.} Interpellato, uno studente indica un
\emph{battle royale}, \emph{Apex Legends}: nel suo ciclo di vita ha più
LiveOps o DevOps? \textbf{LiveOps}. Non che la parte di sviluppo non sia stata
grande, assolutamente; ma l'obiettivo dei \emph{Games as a Service} è
rilasciare il prima possibile e poi iterare, e la loro fase di LiveOps è
estremamente grossa. Ed è difficilissima, probabilmente la cosa più difficile
in assoluto: il prodotto è fuori, non deve crollare a picco, e intanto bisogna
aggiungergli pezzi in corsa. I giochi mobile vivono generalmente di questo, a
meno che l'obiettivo dell'azienda non sia buttare fuori $n$ giochi mobile ``al
secondo'', anche con l'uso dell'intelligenza artificiale, per vedere quale
``si attacca al muro'' e guadagnare un sacco di soldi: una strategia,
purtroppo, dominante.

\paragraph{Il piccolo trucco.} Molte delle cose che sembrano LiveOps in
realtà rientrano nel DevOps. È impensabile, ad esempio, che un DLC venga fatto
interamente dopo il rilascio del gioco: la sua produzione inizia uno o due
mesi prima dell'uscita. Lo si completa dopo, ma gran parte del lavoro viene
portata avanti prima.

\subsection{Qual è la parte più importante?}

Il docente chiede agli studenti quale sia, secondo loro, la parte più
importante del processo, e perché.
\begin{itemize}
  \item Una risposta (Francesca) è la \textbf{post-produzione}, perché è lì
        che ci si occupa di come si vede il gioco.
  \item Un'altra è lo \textbf{start}, l'inizio del progetto.
  \item Federica è indecisa tra \textbf{preproduzione}, perché è la base su
        cui si fonda tutto (e dalle frecce si vede che ci si torna più volte),
        e \textbf{testing}, perché bisogna vedere se ciò che si sta facendo
        funziona.
  \item Un'altra studentessa indica il \textbf{testing}, perché serve a
        mettersi nei panni del giocatore e a rendersi conto delle emozioni che
        il gioco regala.
\end{itemize}

Era una domanda ``a trabocchetto'', come molte delle domande del docente,
ma Francesca e Federica ci sono andate vicino. La fase di rilascio è molto
importante, così come quella che la precede, il \emph{polish} estremo, che è
ciò che rende davvero il gioco piacevole da giocare dal punto di vista del
\emph{feedback}, dei VFX, degli SFX, di tutto il \emph{dressing} che potenzia
le estetiche nel senso di Schell. Anche la preproduzione è estremamente
importante, perché è il fondamento del gioco: se si è creata documentazione e
si è fatta abbastanza ricerca, si hanno fondamenta solide su cui costruire.

Ma la risposta corretta era \textbf{la linea bianca}: la \textbf{freccia che
torna indietro}. È il concetto a cui alludevano Federica e Francesca: il fatto
che si torni spesso indietro. Il ciclo di produzione è \textbf{ricorsivo}. Non
è importante \emph{dove} si ritorni (dalla produzione alla produzione, o
altrove), ma il fatto che si torni indietro, si testi, si veda come stanno
andando le cose e si rivalutino scelte, modifiche e applicazioni. È questa la
parte che rende i giochi (con licenza del termine) ``belli''. Non
necessariamente quelli che vendono di più in assoluto, ma il prodotto mirato,
con i target giusti e le idee che si vogliono trasmettere.

\paragraph{La ricerca nella preproduzione.} Precisando la risposta di
Federica, il docente sottolinea che la preproduzione è anche il momento in cui
si identifica una soluzione a un problema, un buco di mercato, qualcosa che
vende. Esempio: quando sono andati forte i \emph{friendslop}, con
\emph{PEAK} ``e compagnia cantante'', tutti volevano farne, perché costano
meno di un gioco indie ``doppio'' e hanno un ritorno molto elevato: un
friendslop non lo si gioca da soli, lo si gioca con le persone. Cercare la
vendibilità è giusto; ma è altrettanto importante che il prodotto abbia
un'identità e un senso, per creare il miglior prodotto possibile, con un
asterisco: \emph{entro una serie di limiti}, quelli visti nel primo capitolo
(sezione~\ref{sec:limiti}).

\subsection{Tutto quello che facciamo è brutto}

Perché esiste la freccia che torna indietro? Perché \textbf{tutto quello che
si fa, al primo tentativo, è brutto}. Quello che fa il docente, quello che
fanno gli studenti, quello che fa chiunque. Il documento del task, iniziato
presumibilmente uno o due giorni prima, adesso è illeggibile e bruttissimo; lo
sarebbe anche quello del docente, che magari è solo più veloce. La questione
è \textbf{quante volte ci si ritorna}: non ``alla veloce'', ma con senso e con
criterio. Il criterio è vedere come risponde l'utenza, se il livello che si
sta creando è comprensibile, se l'idea implementata evoca le emozioni che ci
si aspetta.

Il lavoro del designer, infatti, è \textbf{far evocare emozioni specifiche nel
giocatore}, e prevedere che cosa il giocatore percepirà giocando è
difficilissimo: impossibile da fare ``a scatola chiusa''. Si può essere il
designer più forte dell'universo, si può essere Miyamoto: ci si andrà vicini,
ma non ci si prenderà mai al 100\%. Ci saranno sempre cose sbagliate: le si
testano, ci si ritorna, le si modificano e si va avanti. Se l'obiettivo è il
gioco migliore possibile entro una serie di limiti, questo processo rende
migliore il gioco, e di conseguenza il designer.

\subsection{Tanti cicli, una sola forma}

Il docente mostra diversi framework di cicli di produzione, da prendere come
linee guida e ``con le pinze'', per far vedere una ricorrenza che torna in
tutti.

\begin{description}
  \item[Il ciclo in tre fasi] \emph{game idea}, \emph{idea refinement} e
        ricerca nella prima parte; la produzione, con una parte a cavallo tra
        le due fasi che comunque continua; la documentazione e il mantenimento
        delle meccaniche, con il loro ciclo di frecce; la post-produzione, con
        il testing dell'intero gioco, il \emph{deployment} (il rilascio) e il
        marketing. Diverse fasi della produzione richiedono \textbf{forze
        lavoro diverse}: i designer servono molto di più nelle prime due fasi;
        i programmatori servono molto in produzione e poi alla fine per
        correggere i bug; gli artisti servono da un certo punto della
        produzione in poi, perché finché le meccaniche vengono validate
        internamente si può ``giocare con i cubi''. Se si dà il prototipo a un
        altro designer dicendogli ``questa cosa dovrebbe essere questo'', va
        bene anche un cubo, un cono, una cosa bruttissima; se lo si deve dare
        in mano al giocatore, va vestito in un altro modo.
  \item[Il ciclo circolare] preproduzione, produzione, testing,
        \emph{pre-launch}, \emph{launch}, \emph{post-launch}: una vera
        circonferenza, sempre ricorsiva.
  \item[Il ciclo dello sviluppo] ideazione, \emph{prototyping},
        \emph{development}, \emph{testing}, \emph{deployment},
        \emph{maintenance}.
  \item[Il ciclo di vita di un prodotto] (\emph{life cycle assessment}) non
        legato ai videogiochi: \emph{manufacturing}, \emph{usage},
        \emph{device recovery}, \emph{service}, \emph{end of life}, pensato
        per il riciclo dei materiali di prodotti digitali (telefoni, computer,
        tablet). Anche qui il ciclo è continuo, e anzi a strati, dal dentro
        al fuori.
  \item[Design, develop, test, deploy, review] vale ad alto e a basso livello.
        Gli \textbf{stakeholder} (altro termine importante), cioè chi detiene
        l'idea o i soldi, come il CEO, compiono questo processo a un livello
        più alto rispetto, ad esempio, al producer che pianifica. Prima si fa
        il design; lo si ``mette a terra'' e lo si sviluppa in piccolo; lo si
        testa al chiuso con una piccola parte di utenza, e poi si amplia piano
        piano (magari non al primo né al secondo, ma al decimo giro); una
        volta sicuri che funzioni, o che non funzioni, lo si implementa nel
        progetto; poi si fa \emph{review}, si controlla che tutto vada bene, si
        aggiusta ciò che non va, si ritesta, e così via.
  \item[Exploring and formulating the problem] esplorazione e formulazione
        del problema; \emph{generating and searching ideas}; \emph{concept};
        \emph{prototyping}; \emph{testing and optimizing}; presentazione. La
        presentazione può valere in diverse fasi: presentare la propria idea,
        anche più volte, ad altri designer, spiegare al lead perché ha senso
        realizzarla rispetto a un'altra; se arriva l'approvazione, secondo un
        processo molto regolamentato, si passa alla fase successiva, che è
        sempre più un discorso di produzione: assicurarsi che i meccanismi
        della macchina girino tutti nello stesso modo.
  \item[Il ciclo in quattro spicchi] analisi (che fa parte della
        preproduzione), sviluppo delle idee, creazione di una soluzione,
        valutazione della soluzione dopo il testing. Con il \emph{design
        brief} come punto di partenza.
\end{description}

\paragraph{Mettere a terra un'idea.} Il docente descrive il passaggio con
più dettaglio. Si ha un'idea, la si è attaccata da una serie di punti di
vista, si è fatta l'analisi e sembra funzionare: si prova a metterla a terra
con un \textbf{prototipo}, qualcosa che sia giocabile in qualsiasi maniera,
anche con le carte, con le scatole di cartone, con i disegni alla lavagna o
con i blueprint di Unreal. Il prototipo \textbf{va giocato}: ``non
dimenticatevi mai quanto è importante giocare le cose''. Anche quando si è
fatto solo un \emph{fix}, bisogna testarlo e controllare che non abbia
distrutto il resto del progetto. Avere un prototipo a sé permette di testare
\textbf{solo quella feature}, in un ambiente chiuso. Ci si mette nei panni
dell'utenza: piacerebbe a questo tipo di giocatore? Risolve il problema? Se la
cosa è perfetta per l'utenza, o almeno si avvicina a ciò che si voleva, la si
mette nel progetto e si passa a un'altra meccanica; altrimenti si migliora e
si ritorna indietro.

\paragraph{Analisi, pianificazione, valutazione.} Nell'ultimo ciclo mostrato
si ritrovano tutte le fasi: l'analisi precedente, ad esempio del mercato e dei
competitor (qualcuno ha già fatto una cosa simile, o proprio la cosa che si
vuole fare?); lo sviluppo e la valutazione delle idee, smontandole o meno; la
\textbf{pianificazione della produzione} dell'idea (quanto costa, quanto tempo
ci vuole, che cosa si può testare prima e che cosa dopo). Si pensi a un
sistema molto complesso come i portali di \emph{Portal}: il primo prototipo
non era perfetto, né testava tutte le meccaniche. Magari c'era un solo
portale, poi due, poi i portali hanno iniziato a modificare la fisica, a
mantenere la velocità del personaggio. Quale di queste cose sviluppare con
priorità e quale dopo? Quale vestire e quale no? Poi si crea la soluzione (un
prototipo, o un piano logico) e la si valuta. Vale per i prototipi e vale per
la documentazione: si fa la prima versione, si controlla, si decide che cosa
modificare, aggiungere o togliere, e si ricontrolla.

\subsection{Il termine chiave: iterativo}

Questi cicli esistono \textbf{per qualsiasi cosa}: per lo sviluppo, la
creazione e l'assemblaggio di auto sportive; per computer, telefoni e
dispositivi elettronici; per lo sviluppo di nuovi farmaci (con la ricerca,
scherza il docente, non sui topi di laboratorio ma ``su grossi blocchi di
tofu''); persino per il design dei vasi da bonsai. Quello che hanno in comune
tutti questi framework, non necessariamente legati al game design, è la
\textbf{forma circolare}: sono ricorsivi, ripetono, ritornano, e c'è un punto
in cui valutano, ordinano, filtrano, per poi tornare al punto precedente.

\begin{quote}
\centering
\textbf{Sono iterativi. Fare design iterativo è l'unico modo per avere un buon
prodotto.}
\end{quote}

Un prodotto migliore del precedente, più vicino a ciò che si vuole fare,
potenzialmente più venduto, più apprezzato dall'utenza, più vicino
all'intento di design di evocare emozioni specifiche nel giocatore: un
prodotto migliore \emph{in toto}, perché a ogni giro si tocca una parte
diversa, o si migliora la stessa. Tutti i prodotti fatti dagli esseri umani,
al primo giro, ``fanno schifo'', e a volte non si possono nemmeno usare.

% ---------------------------------------------------------------------
\section{Iterazione e realtà della produzione: le domande}
\label{sec:domande-iterazione}

\subsection{Se una meccanica strutturale non funziona}

Federica chiede: se alla millesima revisione ci si rende conto che una
meccanica strutturale del gioco non funziona, bisogna cambiare tutto? Se, ad
esempio, sta per uscire \emph{Cyberpunk} e ci si accorge che qualcosa non va,
si cancella tutto o lo si fa uscire comunque?

La risposta ``da mondo del Mulino Bianco'' è: si può anche \textbf{ripartire
da zero}. Il motivo per cui si fa prototipazione è proprio questo:
\textbf{abbandonare velocemente ciò che non funziona} e testare, il più
velocemente possibile, se funzionano altre cose. Se qualcosa non funziona, va
eliminato o aggiornato, e questo può valere anche per il gioco intero.

Ma i videogiochi sono prodotti, con regole specifiche. Una delle ragioni per
cui ci sono fallimenti, o giochi che non ottengono il risultato sperato, è
proprio questa: si è arrivati troppo tardi a modificare o a valutare qualcosa.
Oppure è intervenuto qualcosa di incontrollabile: il mercato si è spostato. Il
docente è sicuro che là fuori ci siano aziende (gliene vengono in mente già
due) che stanno sviluppando un friendslop e potrebbero essere in ritardo:
al rilascio il mercato potrebbe non assorbirlo più, perché saturo, o perché
quel tipo di giocatore si è spostato su altri giochi. Ci sono casi
controllabili e incontrollabili, e il designer deve agire \textbf{su ciò che
può controllare}.

Alcuni punti fermi:
\begin{itemize}
  \item \textbf{Comunicazione e soft skill} (sezione~\ref{sec:soft-skill}):
        non si può andare dal producer o dal capo a dirgli ``sei uno scemo,
        dovevamo farla prima questa cosa'', o ``non possiamo farla perché non
        c'è più tempo''. Bisogna mediare. E spesso è lo stakeholder, o il CEO,
        a dettare a grandi linee che cosa va fatto: si agisce entro quei
        limiti.
  \item \textbf{Limitare i danni} fa parte, purtroppo, del lavoro.
  \item Conoscere il principio dell'iterazione permette di \textbf{ragionare
        già in vista del fatto che le cose dovranno cambiare}. Tanto più
        velocemente si controlla, si osserva, si ``spacchetta'', tanto meno
        si presenterà il problema. Una meccanica non la si ``butta dentro'' il
        progetto appena sembra funzionare: con i \emph{branch}, con altri
        sistemi di versionamento e altre metodologie, ci si assicura prima che
        funzioni e che si sposi con il resto. Non si fa un test, né due, né
        tre: se ne fanno $n$, finché ci sono le risorse.
  \item Il lavoro del designer \textbf{non è buono al primo tentativo} e non
        lo sarà mai (ammesso che il gioco non debba uscire domani). Non si può
        fare \emph{Cyberpunk} in un giorno, ed era chiaro anche agli
        operativi che non si poteva sviluppare \emph{Cyberpunk} in quel tempo:
        lì il problema, senza voler puntare il dito, era principalmente
        produttivo, legato a un sistema specifico adottato. I giochi vanno
        tarati per gestire anche il tempo necessario a riguardare le cose più
        volte.
\end{itemize}
La risposta ``legale'', quindi: se qualcosa non funziona, si cambia, cercando
il miglior compromesso possibile; idealmente, anche distruggendo il gioco.
Ed è per questo che bisogna avere un prototipo il prima possibile: per
testare subito se l'idea ha senso.

\subsection{Se l'ordine arriva dall'alto}

Federica chiede ancora: e se è il produttore a volere a tutti i costi che il
gioco esca, anche se non funziona fin dall'inizio?

Il docente risponde ``da persona operativa, moderatamente frustrata'', perché
gli è capitato non poche volte. Sviluppare progetti è un lavoro di team
(primo capitolo): come il designer si fida del producer, e del fatto che non
gli tiri via il tappeto da sotto i piedi, così il producer si fida del
designer.

\paragraph{A chi si chiedono le stime.} Il docente precisa di non essere, in
quella fase, un lead, un creative director o un game director: è a un livello
più basso. Nelle fasi iniziali la produzione ``a grandi linee'' la si chiede
ai lead o a chi sta sopra. Non si va a chiedere a ogni singolo designer se si
sta scegliendo di fare ``\emph{Cyberpunk 3099}''; non si va dal designer
delle 3C a chiedere quanto ci vorrà per fare la camera del personaggio,
perché bisognerebbe ripetere il processo per 700 persone, il che non è
realistico. Lo si chiede al director, che darà una stima a spanne. Durante la
produzione, di momento in momento, il producer può chiedere direttamente, e
il lead si assicura che i task settimanali o bisettimanali siano nel reame
della fattibilità, ora e in futuro. A quel punto sì, il producer chiede:
``Mattia, secondo te quanto ci vuole per sviluppare questo sistema di
combattimento?''.

\paragraph{Nessuno lo sa davvero.} E qui il docente ammette: \textbf{chi dice
di saperlo è un bugiardo}. Dipende da quanto si è amici del producer, ma la
risposta onesta è ``non lo sappiamo, ti posso dare una stima a spanne''.
Anche quello del producer non è solo un lavoro preventivo: non si può
prevedere esattamente quanto durerà la produzione. Il producer fa prima la
pianificazione e poi il \textbf{damage control}, sistemando i problemi che
sicuramente capiteranno.

\paragraph{Il margine.} Come rendere veritiera una stima? Uno dei ``buoni usi
e costumi'' più semplici è \textbf{aggiungere tempo}. Al designer è capitato
di dare la propria stima aggiungendo tre giorni, e che il producer, benevolo,
ne aggiungesse altri tre: un margine più ampio. Poi può succedere davvero un
disastro: il designer viene investito da un'auto, il progetto ``esplode'',
qualcuno elimina il repository. Lì si tratta di qualcosa di non prevedibile.

\subsection{Publisher, licenze e stakeholder}

Federica cita \emph{Dragon Age: The Veilguard}, chiedendosi come un gioco del
genere riesca alla fine a uscire. Il docente coglie l'occasione per allargare
il discorso. (In una parentesi, cita anche un trailer di \emph{Mass Effect}
realizzato con asset di preproduzione e diffuso da alcuni sviluppatori dopo
la chiusura dello studio da parte di EA: un gesto comprensibile per la
frustrazione, e bello da vedere per capire che cosa il gioco sarebbe voluto
essere, ma \textbf{non professionale}, perché rompe qualsiasi tipo di
NDA.)\footnote{L'episodio è riportato così come raccontato in lezione; non è
stato verificato in questa sede.}

I limiti visti finora (soldi, tempo, numero di persone coinvolte) sono
facilmente accettabili: chiunque deve farci i conti, e possono anche dare una
spinta non indifferente alla creatività. Ma poiché il gioco è un prodotto che
deve essere venduto, entrano in gioco altre questioni. L'azienda può
appoggiarsi a un \textbf{publisher} che richiede una certa meccanica nel
progetto, che vuole dire la sua. Qui prima il director e poi il designer
devono saper \textbf{mediare}, capire che cosa vuole davvero lo stakeholder e
tradurlo in qualcosa che abbia senso nel gioco, mantenendo una certa
\emph{linea editoriale}.

\paragraph{La control room.} Il docente ricollega il discorso a un feedback
ricevuto in mattinata da un'altra classe (Alessandro), da parte di una
\emph{control room}: un ente esterno, che funziona esattamente come un
publisher. Uno studente obietta che, in quel caso, la control room sa bene di
che cosa si parla; vero, ma non è necessariamente così nell'industria. I
publisher potrebbero non sapere nulla di giochi e avere soltanto un enorme
capitale da investire; essendo investitori, vogliono un ritorno altrettanto
ampio. Poiché mettono i soldi, a seconda di ciò che si firma si può
\textbf{cedere parte della propria libertà}. Bisogna saper muoversi in questi
reami: capire che cosa cedere, soprattutto in base a chi chiede qualcosa e a
dove si trova nella gerarchia, per fare il miglior lavoro possibile. Non per
sé (oltre a dormire tranquilli la notte), ma \textbf{per il giocatore}: perché
gli arrivi qualcosa che abbia senso giocare e che gli piaccia, cioè che
evochi il responso emotivo corretto, desiderato. È la parte difficile, un
amalgama di tante cose che si vedranno piano piano.

\paragraph{``Voglio le macchine rosse''.} E se il publisher chiede qualcosa
che non funziona, si è costretti a farlo? Bisogna \textbf{capire che cosa
vuole davvero}. Davanti a un diktat come ``voglio le macchine rosse'', la
prima domanda da porsi è: \emph{perché} le vuole rosse?

\begin{description}
  \item[Le licenze] se si firma un accordo con un marchio automobilistico, il
        marchio vuole che quella macchina finisca nel gioco così com'è: non
        le si possono mettere ``gli alettoni, le ali volanti''. C'è poco
        margine di manovra, e nei giochi simulativi conviene capire quanto
        margine abbiano i designer: spesso la Ferrari vuole che le sue auto
        siano \emph{statisticamente} le più veloci e le più forti del gioco;
        oppure il marchio non vuole che l'auto possa essere distrutta. La
        Ferrari tiene alla propria immagine e non dà licenza per distruggere
        le proprie auto. \emph{Mario Kart} ha collaborato con Mercedes (in
        aula si era ipotizzata la Opel), e lì il problema non si pone: in
        \emph{Mario Kart} le auto non vengono distrutte né modificate; se
        arriva un guscio blu in testa, l'auto resta uguale.\footnote{Si tratta
        del contenuto aggiuntivo gratuito Mercedes-Benz per \emph{Mario Kart
        8} (2014).}
  \item[Il caso Wipeout] in un gioco di corse (il docente cita
        \emph{Wipeout}, pur non essendone certo), il publisher lamentava che
        le macchine fossero troppo lente. I designer erano sicuri che il
        sistema funzionasse correttamente, e il feedback del publisher non era
        davvero ``le macchine vanno troppo lente'', ma ``fatemele più
        veloci''. Il compito del designer è \textbf{continuare a fare
        domande} fino a trovare la \textbf{radice del problema}: le macchine
        non venivano \emph{percepite} come veloci. Due effetti a schermo, due
        effetti di lente della camera, due VFX, due scie: problema risolto
        senza toccare nessuna statistica. Anche questo è lavoro di design.
\end{description}

\paragraph{Questioni che non riguardano il design.} Ci sono poi questioni
economiche: quanto costano i team, o perché certe acquisizioni vengono fatte
(ad esempio per i ritorni e le agevolazioni fiscali legate alla quantità di
persone e di team che si assorbono). Sono questioni complesse, di
mantenimento dell'impresa e di marketing: riguardano il design solo se si è
CEO, e non verranno trattate, probabilmente nemmeno nei tre anni, a meno di un
workshop apposito.

% ---------------------------------------------------------------------
\section{I principi dell'iterazione}
\label{sec:principi-iterazione}

\begin{quote}
\centering
\textbf{L'unico modo per rendere un prodotto sempre migliore, sempre meno
peggiore, è il processo iterativo.}
\end{quote}

È \textbf{la prima regola} da ``stamparsi in testa'': assorbirla, masticarla,
ingoiarla completamente. I progetti brutti, terribili (con un'arte
terrificante, a cui manca un pezzo, pieni di bug), spesso lo sono perché sono
passati poche volte attraverso questo processo. È il primo metro di misura.

Uno studente fa notare che vale anche per le \textbf{animazioni}. Certo: in
un sistema di combattimento le animazioni iniziali sono provvisorie. Si può
trovare in rete il prototipo del combattimento di \emph{Bayonetta}: ``fa
schifo'', ci sono figure placeholder e pose provvisorie che si alternano ogni
due secondi, ma dà al designer l'idea delle tipologie di animazione legate al
combattimento, come il \emph{wind-up} e l'\emph{anticipation}. Questo
discorso si fa (o si dovrebbe fare) in tutti i dipartimenti.

Dal principio dell'iterazione discendono alcuni corollari.

\subsection{Target audience}

I prodotti sono creati per uno specifico \textbf{target audience}, abbreviato
\textbf{TA}. L'obiettivo del professionista è creare \emph{per altri}: il
sudore, il tempo e il lavoro vanno fatti in direzione di quei giocatori.

\paragraph{E Balatro?} Una domanda (che il docente attribuisce a Daniele, e
che dice gli abbia ``aperto il mondo'') è: perché allora lo sviluppatore di
\emph{Balatro} ha avuto tanto successo facendo, di fatto, il gioco per sé? La
risposta è semplicistica, e ci sono molte altre questioni da valutare, ma la
prima, per il discorso sul target audience, è che \textbf{era lui il proprio
target audience}, e ``per puro caso'' di persone come lui, con necessità e
richieste simili, ce ne sono parecchie. Vale per molti altri sviluppatori
indipendenti: Toby Fox con \emph{Undertale}, per esempio. Poi si possono
investigare altri motivi del successo: \emph{Undertale} è pieno di citazioni e
di riferimenti al mondo di internet e a un certo tipo di cultura moderna;
\emph{Stardew Valley} ha ampliato molti sistemi di \emph{Harvest Moon} (il
ciclo giorno--notte, le relazioni con gli NPC, la possibilità di
personalizzare la propria fattoria, la progressione per cui si sblocca prima
una zona, poi l'isola, poi il vulcano). Il docente invita quindi a
``parafrasare'' con cautela il caso dei giochi fatti per sé: \textbf{ogni gioco
ha un target audience}, e va capito quanto è ampio o stretto.

\paragraph{Soddisfare l'utenza.} Bisogna prestare attenzione all'utenza,
cercare le cose che le piacciono, che cerca, o che la stimolano nella
direzione corretta. La madre del docente non giocherebbe mai un gioco che
punta a essere estremamente eccitante, o che fa paura, né un MMO: interpellata
in aula (con i titoli che gioca), si capisce che non cerca il \emph{grinding}
estremo o la socialità estrema di un gruppo di cinquanta persone che fanno
\emph{raid}. Gli elementi implementati nel gioco sono per il target audience.
Ci saranno compromessi con gli altri dipartimenti, tutto deve ``cliccare''
insieme, ma l'obiettivo finale è questo:

\begin{quote}
\centering
\textbf{Ogni decisione di design è fatta per soddisfare il target audience.}
\end{quote}

Ne consegue che, quando si discute con un collega, non lo si fa per avere
ragione o per dimostrare che il proprio design è migliore: si sceglie
l'opzione \textbf{più sensata per un terzo}, l'audience. Non ``ho ragione io'',
non ``voglio questa parte, data da me, nel gioco''. Il gioco è \emph{di tutti
e di nessuno} di quelli che ci lavorano: è di chi lo riceve.

\subsection{Involvement}

Il target audience dovrebbe essere \textbf{coinvolto} in parte della
produzione; va capito quanto. Nei giochi indipendenti l'audience è chiamata
spesso in causa, anche per necessità e per ragioni di marketing; nei giochi
più grandi non la si coinvolge nei prototipi molto brutti, ma esistono gruppi
di testing molto ridotti, composti da giocatori che potrebbero essere il
target ideale, che testano il gioco e danno feedback.

\begin{itemize}
  \item Va testato soprattutto il \textbf{core concept}, l'idea iniziale, per
        capire se ha senso. Un Metroidvania con dinosauri fatati, per il
        docente, è un'idea fantastica; ma magari solo per lui, e nessuno là
        fuori ha bisogno di quel gioco: allora forse è il caso di non farlo.
        Magari i dinosauri li vogliono demoniaci, e allora lo comprerebbero
        tutti. Forse: servono i dati, serve la ricerca.
  \item Vanno valutate \textbf{build testabili}, per assicurarsi che il gioco
        venga compreso. Capita continuamente di creare un livello in cui per
        il designer è tutto chiarissimo (la progressione, dove andare, che
        cosa fare) e di darlo in mano a qualcuno che non l'ha mai giocato e che
        non ha idea di che cosa fare.
  \item Quando si dà una build a un \textbf{publisher}, deve essere ``a prova
        di scemo'': molti publisher non sono giocatori, e lo stesso vale per
        alcuni giornalisti. Il docente cita il famoso video di un giornalista
        alle prese con la demo di \emph{Cuphead}, incapace di superare il
        tutorial (bisognava premere un salto e poi uno scatto in aria, cioè
        due tasti in sequenza).\footnote{Il video è quello registrato da Dean
        Takahashi di VentureBeat alla Gamescom 2017, diventato virale proprio
        per la difficoltà mostrata nel tutorial.} ``Il giornalismo videoludico
        ha bisogno di un'evoluzione'', commenta; ma se il gioco lo si deve
        vendere alla propria madre, non le si può mettere un ostacolo che
        richiede di premere più di un tasto alla volta.
  \item Si osserva anche, lato marketing, come risponde l'utenza: che cosa
        piace e che cosa no. I \emph{reel} ``brain rot'' funzionano, ahimè;
        e tutti conoscono le pubblicità dei giochi mobile con gli omini che
        sparano e poi falliscono: sono fatte così per una ragione precisa,
        perché attirano la stragrande maggioranza dei giocatori.
\end{itemize}

\subsection{Iteration}

I giochi evolvono e migliorano grazie al \textbf{feedback} di utenti ed
esperti. Prima si fa testare il proprio design, il proprio prototipo, a un
altro designer, possibilmente uno che conosce poco il progetto (meno lo
conosce, meglio è), e si verifica se recepisce le informazioni nella maniera
corretta: se capisce dove deve andare, che cosa deve fare, se la storia che
recepisce è quella che gli si vuole raccontare. Gli studenti di narrativa più
tendenti al \emph{worldbuilding} che al \emph{narrative design}, ad esempio,
tendono a scrivere storie estremamente complesse, con ``27 milioni di nomi''
e neologismi impossibili da ricordare (``il Bibus legato al re Bobus, del
regno Namus'').

Bisogna quindi testare, controllare e \textbf{chiedere}: hai capito questa
cosa? Come l'hai recepita? Ti è piaciuta? Perché? Che emozioni ti ha dato?
Una volta interiorizzata la lista delle emozioni
(sezione~\ref{sec:vocabolario}), si può anche \textbf{ridurre il passo}:
fare domande volutamente informali, che non si farebbero a un altro designer
(``ti piace sparare? ti piace esplorare? ti piace parlare con altre persone?
ti piace sbloccare achievement?''), e ricevere in cambio molte più
informazioni. Si prendono le emozioni, le si accorpano in gruppi e si cerca
la domanda più generica possibile che però indirizzi nel modo migliore verso
uno di quei gruppi. HoYoverse (ex miHoYo, gli sviluppatori di \emph{Genshin
Impact}) è estremamente brava a fare questo tipo di sondaggi, e il docente
ammette di ``rubare'' molto dalle loro domande. Forse, se ci sarà tempo, si
farà un ``Survey 101''.

\subsection{Holistic}

Sempre per ragioni di produzione, l'idea è di \textbf{sviluppare in parallelo}
quante più cose possibili, entro dei limiti.

\paragraph{Il block out.} Il docente chiede quanti sappiano che cos'è un
\emph{block out}; è giusto non saperlo. I livelli si fanno così: il
\textbf{level designer} lavora sugli \textbf{ingombri geometrici}, dà l'idea
del livello e la documenta (se il blocco di ``tofu'' rappresenta una magione
romana, bisogna spiegare come dovrebbe essere e che cosa contiene); poi
l'\textbf{environment artist} lo \emph{veste}, lo ``dressa'', mettendoci sopra
i \emph{prop}. È come dare a un bambino un disegno da colorare: il level
designer traccia i contorni, gli artisti ci colorano dentro. Solo che, nel
50\% dei casi (forse meno), \textbf{gli ingombri non vengono rispettati}: un
albero si pianta davanti al giocatore e non lo lascia proseguire. Bisogna
quindi testare il livello e modificarlo: è il discorso dell'iterazione, e
succede sempre, tanto più quante più persone lavorano al progetto.

Se invece si sviluppa insieme (mentre si fa il prototipo della meccanica, il
VFX artist fa l'effetto, l'SFX artist fa il suono, il programmatore inizia a
montare le cose e il designer osserva che tutto vada in linea, che i tempi
siano rispettati), si va più veloci: si raggiunge prima qualcosa di giocabile,
lo si testa prima, lo si itera prima e più volte, e il gioco, per
un'equazione semplice, diventa migliore.

\paragraph{Che cosa significa olistico.} Uno studente propone una definizione:
\emph{olistico} è ciò che raccoglie più questioni; secondo una visione
olistica, $x + x + x$ non è uguale a $3x$, ma a qualcosa di maggiore, perché la
visione è l'insieme di tutto. Il docente precisa che qui il senso è più
pratico: \textbf{includere tutti i dipartimenti}, più sfere possibili. Se si
porta il proprio design (la meccanica, il sistema) il prima possibile davanti
agli altri dipartimenti, invece di rimbalzare tra team diversi con un'idea
già terminata, il programmatore, l'artista o il sound designer possono vedere
qualcosa che il designer non vede, o un limite del loro dipartimento. L'idea
può venire da chiunque, ma ognuno è specializzato nel proprio campo: se
l'artista dice ``questa struttura modulare non riusciamo a farla, dobbiamo
farla così'', si può già pensare a un'altra soluzione.

C'è però un ordine: la pipeline ha dei passaggi precedenti e successivi.
L'artista non si mette a vestire il livello se il livello non c'è, né mentre
il level designer lo sta costruendo: prima lo si costruisce, lo si consegna
all'artista, che lo veste e lo restituisce, e lo si testa. Allo stesso modo,
prima si fa il design e poi il programmatore lo implementa, altrimenti si
inventano cose che vanno rifatte. Ci si può portare avanti, ma l'obiettivo è
\textbf{incastrare i lavori nel modo più efficiente possibile}: è compito del
producer, e fa parte dell'ottimizzazione della macchina che permette
iterazioni più veloci e più numerose.

\subsection{Adaptive}

È il discorso fatto con Federica e a proposito del producer: \textbf{il gioco
cambia}, per forza. Bisogna essere bravi a prendere la palla al balzo e
trovare il cambiamento più efficace.

Ci sono però cose che dovrebbero restare quelle: i \textbf{pilastri} del
gioco, che si vedranno in un'altra lezione di documentazione, legata a un
documento specifico. Quelli non si cambiano. Se vengono meno, \textbf{non è
più lo stesso gioco}, non è la stessa esperienza, e a quel punto si fa
un'altra cosa. È anche il motivo per cui spesso le produzioni sono
infinitamente lunghe, o arrivano al traguardo con qualcosa di incoerente o
molto diverso da ciò per cui si era fatta preproduzione: non si sono
mantenuti, salvaguardati e condivisi i pilastri all'interno del team. Se uno
di essi viene meno o non è condiviso, il progetto non c'è più, ed è molto più
vantaggioso ``tirare via tutto'' e ricominciare da capo.

\subsection{I segreti per iterare}
\label{sec:segreti-iterazione}

\paragraph{Create un prototipo il prima possibile.} Dopo la fase di ricerca
e di documentazione, il prototipo va creato \textbf{il prima possibile}.
\begin{itemize}
  \item Deve \textbf{costare poco}: non ci si deve spendere troppo tempo.
  \item Deve essere \textbf{efficiente} nel testare ciò che si vuole testare.
  \item Deve essere \textbf{abbastanza completo} rispetto a ciò che si deve
        testare: una feature (un sistema di camere, un sistema di
        locomozione) o l'intero progetto.
  \item Deve essere \textbf{giocabile}, non importa come: digitale, gioco di
        carte, scatoloni, o il docente e Valerio che prendono una spada di
        cartone e si ``menano''. Qualsiasi cosa può funzionare: un LARP
        (\emph{live action role-playing}), cioè persone in un prato con le
        asce di polistirolo; D\&D con le persone; persino il \emph{modding} di
        altri giochi.
\end{itemize}

\paragraph{Il prototipo serve anche a comunicare.} Uno studente chiede se
anche per una presentazione del concept del progetto convenga avere già una
build. Sì: è utile non solo per validare internamente se il gioco funziona, ma
anche per \textbf{far capire a un altro che gioco si intende}. Non si ha idea
di quanto sia difficile comunicare un gioco che non esiste. Anche i
\emph{core elements} della lezione precedente servono a questo: partire da dei
pezzi, come le \emph{game qualities} di Romero e Schreiber, dà già un'idea.
``Metroidvania'' può essere qualsiasi cosa; ``Metroidvania gotico, ambientato
in un mondo di insetti'' (\dots e magari ``se ne potrebbe pure fare un
sequel'') fa già capire a tutti che si parla di \emph{Hollow Knight}. Più
elementi chiari si aggiungono, più il gioco è chiaro. E se il gioco lo si può
già giocare (se si vede il cavaliere e lo si può muovere, se si vede una
concept art di come si presenta il mondo, o anche solo un quadrato con uno
smile che muove un ago contro un blocco rosso, e si sente il \emph{feeling}
del salto, del colpo, della parata), il gioco \textbf{prende vita}, e lo si
può comunicare, testare e valutare.

\paragraph{Fate giocare i prototipi.} I prototipi non restano ``lì nel
nulla'': vanno giocati da esperti e da parte del target audience, con dei
limiti. Se la build è lacunosa in alcuni punti, lo si dice al giocatore e si
creano dei \emph{limiti} (``non date feedback su questa parte, concentratevi
su quest'altra''), così come il docente aveva chiesto alla classe di non
badare all'aspetto delle proprie slide. Si prendono appunti e dati, non
giudizi del tipo bello/brutto: si crea un form, si fanno domande, si cerca di
ottenere quante più informazioni possibili. Poi si \textbf{analizzano} i dati
ottenuti, per migliorare in piccolo quel design o quella feature, o in grande
il progetto. E lo si fa non una, né due, né tre volte, ma \textbf{finché si ha
il tempo di farlo}.

\paragraph{Il primo segreto.} Il motivo per cui, soprattutto in passato, i
tripla A erano così gradevoli da giocare è che questo processo iterativo
veniva fatto per anni; potevano permetterselo perché avevano molti soldi. Ma
se le risorse vengono sfruttate così, il guadagno, quello \emph{artistico}, è
assicurato.
\begin{quote}
\centering
\textbf{Più è lungo il processo iterativo, migliore è il risultato.} Prima ci
si arriva, prima lo si fa, più lo si fa, e meglio il gioco diventa.
\end{quote}

\paragraph{Il ciclo rivisto.} Ecco perché il primo ciclo mostrato era
bruttissimo: anche quello si può migliorare, si può iterare. Il docente lo
ripresenta arricchito:
\begin{itemize}
  \item la freccia dell'iterazione, come suggeriva Federica, può arrivare fino
        alla preproduzione, ma l'iterazione si fa sia sull'intero gioco sia,
        con frequenza ancora maggiore, sulle singole feature in produzione;
  \item non c'è solo la beta: ci sono diverse \textbf{milestone} e diverse
        \textbf{build}, le più note delle quali sono \emph{alpha}, \emph{beta}
        e \emph{gold} (il loro significato verrà spiegato più avanti);
  \item dopo il rilascio si può proseguire con le LiveOps;
  \item tra l'inizio del progetto e la preproduzione c'è la \textbf{visione}:
        fa parte della produzione, ma è il momento in cui si condivide con il
        team la visione e i famosi pilastri. Se la preproduzione è lo
        scheletro della casa, la visione è la \emph{fondamenta} del progetto,
        ``la pavimentazione della pavimentazione''; i pilastri sono i massicci
        blocchi di cemento che tengono su la struttura;
  \item in \textbf{preproduzione} si fanno generalmente visione del progetto,
        ricerca, analisi, prototipazione e documentazione;
  \item in \textbf{produzione} si fanno implementazione, \emph{dressing},
        \emph{management}, \emph{tuning} (il bilanciamento, di cui si è
        parlato nel primo capitolo) e, ancora, documentazione, oltre alla
        prototipazione, che continua.
\end{itemize}
Il disegno è brutto perché l'idea è di iterare: prima lo si fa su Paint, in
modo terribile, ``con i cosi sbagliati'' (anche per comodità, soprattutto per
le modifiche), poi ci si pensa a migliorarlo. Non sarà mai la versione
finale.

\paragraph{Il secondo segreto.} ``Vi ho mentito'', ammette il docente: i
segreti sono due.
\begin{quote}
\centering
\textbf{Prima createlo, poi fatelo diventare bello.}
\end{quote}
Bisogna creare qualcosa di brutto, e non lasciare che il documento o il task
resti una pagina bianca per paura di buttarci sopra delle cose. Si creano
dati, si guardano video, si scrive, si butta giù qualsiasi cosa; poi si
taglia, si ``trincia''. Queste due regole servono nel lavoro di tutti i
giorni, adesso, e serviranno a breve, quando si parlerà di progetti.

% ---------------------------------------------------------------------
\section{Perché la documentazione è importante}
\label{sec:perche-documentare}

\subsection{Una pausa e un promemoria}

Dopo una pausa di sette minuti, il docente ricorda di aver dimenticato di
parlare di un concetto di \textbf{Raph Koster}, che non tratta subito per
paura che prenda troppo tempo: chiede agli studenti di ricordarglielo nella
lezione di lunedì. Si tratta di qualcosa di molto utile e introduttivo,
legato alla lezione su che cos'è un gioco (capitolo~\ref{cap:che-cos-e-un-gioco}):
non proprio un framework (anche se un framework c'è), ma un concetto da cui
sono stati derivati altri framework, legato alla mentalità di \emph{dove sta
il gioco nel play}, cioè che ruolo ha il giocare e soprattutto come viene
\textbf{percepito}. Era emerso in alcune definizioni degli studenti, e Koster
lo spiega in maniera regolata e limitata.

\paragraph{Il formato del task.} Una studentessa chiede se il documento del
task debba essere un testo libero o una serie di paragrafi. Il docente ne ha
parlato a inizio lezione (sezione~\ref{sec:drive-regole}) e invita a
riguardarla: per ora valgono soltanto le regole scritte nel documento delle
regole dei task, più le precisazioni date a voce.

\subsection{La prima vera lezione di documentazione}

Dopo la lezione zero, inizia la prima vera lezione di documentazione. Anche
questa si apre con un meme, ``estremamente vero'': se un professionista ti dà
un documento terribile, \textbf{è una persona cattiva}.

Perché la documentazione è così importante, e perché il docente la
sottolinea in modo particolare (oltre al fatto che è una sua specialità)?
\begin{itemize}
  \item Perché \textbf{è il primo lavoro del designer}, la prima cosa che
        deve fare.
  \item Perché la capacità di documentare e di \textbf{adattarsi a diversi
        formati} di documentazione determina la differenza tra un lead e un
        junior o un intern, tra un game director e gli altri. Non in toto, ma
        la documentazione è già una buona \emph{cartina tornasole}
        dell'esperienza di un designer.
\end{itemize}

\paragraph{Genji e Mercy.} Nella slide compare Genji, il personaggio di
\emph{Overwatch}: il DPS che va in giro a dettare ordini e a ``sparare
documenti in testa alle persone''. Il designer \textbf{non è} questo. Il
designer è \textbf{Mercy}, il \emph{support}: la persona che tende la mano agli
altri.

\subsection{A che cosa serve la documentazione}

\begin{description}
  \item[Mantenere la visione coerente] tra tutti i membri del team, ai
        rispettivi livelli. È il primo punto: i pilastri di cemento devono
        essere visti allo stesso modo da tutti. Non deve succedere che ognuno
        veda una cosa diversa: che lo si guardi dal dipartimento di arte o da
        quello di design, il pilastro è lo stesso, la visione è la stessa. E
        l'unico modo per ottenerlo è \textbf{comunicare coerentemente}.
  \item[Aggiornare il lettore] chi non sa nulla ed entra nel team, leggendo la
        documentazione, si mette in pari (\emph{up to date}) senza dover
        parlare con tutti. Parlare è utile, ma non tutti avranno tempo per
        raccontare il gioco per una giornata intera, soprattutto nelle aziende
        molto grandi. I documenti raccontano la storia del team, delle persone
        e soprattutto di quale dovrebbe essere il progetto, di qual è la
        direzione.
  \item[Ricerca] una buona documentazione rende i dati non ``noiosi'': le
        persone leggeranno i dati e la documentazione con piacere, senza
        esserne spaventate. La soglia di attenzione è stata ``completamente
        annichilita'': quanti, aprendo un libro o un documento, guardano per
        prima cosa quanto è grande la barra di scorrimento, per prepararsi a
        capire quanto è ``gargantuesco''? Le cose che attirano sono quelle ben
        documentate. Per questo si cerca di dare un focus specifico a un certo
        dato, o di metterlo in grafico, perché sia comprensibile e facilmente
        digeribile. Il lato oscuro di questa capacità è la
        \textbf{manipolazione dei dati}: ``da grandi poteri derivano grandi
        responsabilità''. Va usata onestamente, perché l'obiettivo non è far
        valere la propria idea, ma trovare l'idea migliore per il target
        audience.
  \item[Supporto] è il motivo per cui c'è Mercy: aiutare gli altri. Con la
        documentazione si aiutano gli altri dipartimenti a ragionare come il
        designer, a capire ciò che ha pensato e poi a metterlo a terra.
        Invece di passare tutto il tempo ``sulla spalla'' del programmatore a
        dirgli fai questo, fai quello, no così non va, gli si dà un documento
        con esattamente ciò che si vuole, nel modo esatto in cui lo si vuole,
        e il programmatore si preoccuperà di implementarlo. Vale anche per gli
        artisti: dove è ambientato il livello, come dovrebbe essere vestito.
        In un gioco ambientato a Babilonia: come ragionavano i babilonesi,
        com'erano strutturate le strade, com'era fatta una ziggurat. Così gli
        altri vedono un mondo che ancora non esiste nello stesso modo in cui
        lo vede il designer. Per questo il designer deve avere la visione il
        più chiara possibile: se ce l'ha chiara e la documenta nel modo
        migliore, \textbf{anche gli altri vedono ciò che non si vede}.
  \item[Controllo] la documentazione permette di tenere sotto controllo la
        produzione: spacchettare le cose che servono, averle localizzate (sul
        Drive, su Confluence, su Miro, dove si vuole). Se si sa quali passi
        fare e quando, si sa anche se ci sono problemi: se ci si sta arenando
        su una parte di una feature, o non si riesce a raggiungere una certa
        versione, lo si può affrontare con il team e con il producer. La
        documentazione serve a tracciare lo stato del progetto, nel grande e
        nelle sue piccole parti. Nessun designer documenta l'intero progetto:
        ognuno si occupa di una parte, e insieme fanno il tutto.
\end{description}

\subsection{I documenti evolvono}

\begin{quote}
\centering
\textbf{I documenti evolvono con il team e con il progetto.}
\end{quote}
I documenti non restano scolpiti nella pietra, tranne alcuni. Il
\textbf{documento di visione}, ad esempio, solidifica i pilastri: si fa e
così resta. Ma gli altri, come il \emph{game design document} di una feature,
vengono aggiornati sulla base del processo iterativo: alla data del 23
ottobre la feature è così; si testa, si itera, non va bene, si aggiorna il
documento; il 25 o 26 ottobre ci sarà un nuovo documento con le informazioni
sullo stato della produzione in quel momento esatto.

\subsection{La vostra firma e la vostra arte}

\paragraph{Il biglietto da visita.} La documentazione è la prima cosa che si
fa ed è la propria \textbf{firma}, il proprio biglietto da visita, anche per
il portfolio. I \textbf{test di selezione} delle aziende (Ubisoft, CD Projekt
Red, Stormind e così via) sono generalmente \textbf{documenti}. Per un ruolo
da level designer probabilmente verranno chiesti anche dei block out; per un
ruolo da technical designer una bozza o un prototipo, ad esempio di un
combattimento. Ma per il game designer (e a volte, è capitato, anche per il
technical) il documento è sufficiente: chiedono proprio quello.

\paragraph{La vostra arte.} Il docente sposta il discorso su un piano più
romantico, che gli sta particolarmente a cuore: per lui, ``Mattia Podini
designer'', fare documentazione è \textbf{piacevole}. Potrebbe non esserlo
per tutti, ma se si trova il proprio \emph{sweet spot}, un modo di ragionare e
di documentare piacevole per sé in primo luogo, il lavoro diventa molto più
semplice e anche più apprezzato. Non a tutti piace, ma fa parte del lavoro:
trovare un modo buono e piacevole di mettere a terra e condividere le idee, di
scrivere. ``Questa roba la dovete fare: vi pagheranno per farla.''

% ---------------------------------------------------------------------
\section{Le golden rules della documentazione}
\label{sec:golden-rules}

Esistono delle regole, le \textbf{golden rules}, per scrivere buona
documentazione. Non sono dogmi scolpiti nella pietra: il docente stesso ne
rompe alcune. Ma lo fa perché ha anni di esperienza alle spalle e conosce
``un paio di segreti'' della documentazione; le conosce tutte perché le ha
\emph{subite} sulla propria pelle quando studiava, perché è stato limitato
tante volte, e ormai ne conosce il valore. Gli studenti no: quindi
\textbf{sotto queste regole ci si sta e non le si rompe}, pena la non lettura.
(Sono ``i dieci comandamenti'', anche se in realtà sono un po' di più.)

\subsection{Mostra, non scrivere}

Il modo migliore per comunicare qualcosa è \textbf{farlo vedere}, se si può:
immagini, video, \emph{flowchart} (si vedrà che cosa sono), link. Le
informazioni vanno condensate in maniera visiva. Se si sta facendo il design
di un menu, se si spiega una meccanica, si fa vedere l'immagine: il
protagonista che attacca un nemico. Anche uno \emph{stickman} con lo spadone e
un nemico triste va benissimo (``punto uno, punto due: gli stacco la testa,
esce un sacco di sangue, lui a terra''): un'immagine da sola non basta, ne
servono di più, ma è più facile che una persona capisca con le immagini che
con le parole. Una cosa non esclude l'altra: si può scrivere il testo e
affiancargli l'immagine.

\paragraph{Mock-up e wireframe.} Se si parla di UI si creano dei
\emph{mock-up}, i cosiddetti \textbf{wireframe}. Devono essere brutti (o
meglio: più belli si riescono a fare, meglio è; ma averli brutti ha due
vantaggi). Il docente chiede alla classe quali siano, insistendo perché gli
studenti parlino e facciano domande (``c'è il solito discorso sulla
comunicazione''):
\begin{enumerate}
  \item \textbf{Ci vuole meno tempo}, e quindi, per usare la parola giusta
        (suggerita da uno studente), \textbf{si può iterare più
        velocemente}. Ed è comunque molto chiaro: si vedono gli spazi, come
        sono organizzati, dove devono andare le cose.
  \item Il secondo vantaggio è ``politico'': si \textbf{evita di pestare i
        piedi} a un altro dipartimento.
\end{enumerate}

\paragraph{Esempio: Hollow Knight e Zelda.} Si supponga di avere un sistema
simile a quello dei \emph{charm} di \emph{Hollow Knight}. Si confrontano due
menu: quello di \emph{Hollow Knight} e quello di uno \emph{Zelda}. Entrambi
sono giochi con combattimento corpo a corpo a breve distanza, con un'arma
principale (l'ago in \emph{Hollow Knight}, la Master Sword in \emph{Zelda});
\emph{Hollow Knight} ha un sistema di magie che \emph{Zelda} non ha, mentre
\emph{Zelda} ha bombe, arco e così via; in \emph{Hollow Knight} ci sono le
maschere che aumentano la vita, in \emph{Zelda} i cuori; ci sono gli oggetti
(come le chiavi), e in \emph{Hollow Knight} c'è una descrizione per ogni
oggetto, mentre in \emph{Zelda} no. Sono, fondamentalmente, molto simili: menu
di inventario. Ma se il designer, invece di fare un wireframe neutro, colora il
menu di arancione e al posto dell'ago mette una spada draconica trovata
cercando ``sword PNG'', sta già comunicando qualcosa sullo \textbf{stile
artistico}, che non è di sua competenza né del suo dipartimento: non è un
direttore artistico né un UI designer. Se gli viene chiesto di fare la
\textbf{struttura} del menu e le sue \textbf{interazioni}, il suo compito è
fare in modo che premendo un tasto si vada sui charm, premendone un altro si
vada sulla mappa, che si veda o meno la descrizione (come in \emph{Zelda}):
queste cose le decide lui. (Esiste una branca del design dedicata, quella
dell'UX/UI design.) Se non si ha l'\emph{ownership} artistica, vestire il
wireframe può essere deleterio, perché comunica qualcosa che non dovrebbe.

Le alternative sono due: scrivere ``a caratteri cubitali'' sopra le immagini
che si tratta di \emph{render artistici non definitivi}; oppure, come fa il
docente con la narrativa, dare comunque dei nomi (invece di ``Pinco Pallino''
o ``nome TBD''), segnalarli in un certo modo e mettere prima un
\textbf{disclaimer}: tutte le definizioni narrative sono da rivedere,
sentitevi liberi di prenderle, smontarle e distruggerle. Ma va detto.

Un wireframe permette comunque di capire quanto dovrebbe essere grande un
elemento nello spazio, se è una schermata intera, se dietro c'è la schermata
di gioco o no, dove si va premendo un tasto, se bisogna tenerlo premuto per
andare altrove, e così via. Se lo si fa vedere, in generale, si capisce di più.

\subsection{Essenziale}

I documenti devono essere \textbf{essenziali}. Non si scrivono documenti di
50 pagine, perché nessuno li legge: le persone che dovrebbero leggerli non
vogliono farlo, e se non vogliono non lo fanno, anche se è il loro lavoro;
oppure li leggono controvoglia e non capiscono, o non ricordano, nulla.
\begin{itemize}
  \item Concentrarsi sull'\textbf{obiettivo} del documento, su ciò che si
        deve comunicare.
  \item Essere \textbf{chiari e concisi}.
  \item Non usare \textbf{parole deboli} come ``forse sarà così'', ``potrebbe
        essere questo o quello''. Se un dato è quello, si scrive il dato: se
        il personaggio ha 50 punti vita, si scrive 50, non ``tra i 5 e i
        100''. Se non si ha un'idea precisa, si scrive comunque un valore, ad
        esempio ``valore intero tra 0 e 100'', e poi si vedrà.
\end{itemize}

Attenzione: essere essenziali non significa fare documenti di una pagina
senza pensare attentamente a che cosa scriverci. È \textbf{più difficile} che
scrivere 70 pagine, perché quelle 70 pagine magari si scrivono comunque, ma
poi si procede \textbf{per sottrazione}: si parte da tante pagine e si taglia
ciò che non serve. È anch'esso un processo iterativo: si parte dal grosso e si
arriva al piccolo.

\subsection{The right tool for the right job}

Bisogna scegliere il \textbf{supporto ideale} per ciò che si deve comunicare
(una domanda che aveva posto anche Antonino). Se si deve presentare un'idea,
fare un \emph{pitch}, si usano le \textbf{slide}: PowerPoint, Google Slides,
Canva o qualsiasi strumento, purché permetta un largo uso di immagini. Se si
deve bilanciare qualcosa, non lo si fa su Google Docs o su Word, ma con un
foglio di calcolo (Excel o Google Sheets). Vale per il supporto fisico e vale
per il contenuto, esattamente come per i framework: si sceglie il framework
giusto per il gioco giusto, il tempo giusto, le risorse giuste. Raramente ce
n'è uno migliore di tutti; ce ne sono alcuni migliori in determinati contesti
(si veda la sezione~\ref{sec:scegliere-documento}).

\subsection{Coerenza e sequenza}

Bisogna comunicare in maniera \textbf{coerente} e con la \textbf{sequenza
corretta}. Non si parla di una meccanica, poi si salta a un'altra, poi si va
altrove e poi si rimanda indietro; né si scrive qualcosa a cui manca
un'informazione che compare cinquanta paragrafi sotto.

Rientra in questo punto anche la \textbf{lingua}: niente errori grammaticali
gravi. Eventuali difficoltà specifiche (dislessia, discalculia e simili, di
cui si era parlato nel primo capitolo) vengono ovviamente considerate, ma
bisogna cercare di essere compresi e coerenti: tempi verbali giusti, rapporto
tra coordinate e subordinate corretto, e così via.

\subsection{Mind the language}

Se esiste una \textbf{nomenclatura tecnica}, la si sfrutta. Così come non si
possono usare le parole bandite, si devono usare le parole e i framework visti
a lezione. Come aveva detto il direttore: se il linguaggio condiviso è quello
tecnico, e si parla da professionisti, si viene riconosciuti come tali.

\subsection{Niente decorazioni inutili}

Fondamentalmente: \textbf{evitare la ``monnezza''}.
\begin{itemize}
  \item Niente decorazioni per abbellire il documento: se un elemento non è
        utile a comunicare qualcosa, va tolto.
  \item Niente pagine (quasi) vuote: se in un documento ci sono intere pagine
        vuote, il docente non lo legge. I fogli A4 ``sono oro'': lo spazio
        disponibile va sfruttato tutto, ``come l'acqua nel deserto''.
        Bisogna pensare a ogni spazio e a ogni parola che si scrive.
  \item Non si sta scrivendo una tesi: il linguaggio non è quello accademico
        in senso lato. I documenti tecnici non sono testi fitti, e hanno regole
        specifiche, diverse da quelle accademiche.
\end{itemize}

\subsection{Grafici con intento}

Quando si può, i dati inseriti nei documenti vanno messi in \textbf{grafico},
in linea con il primo punto (mostra, non scrivere), ma con \textbf{cognizione
di causa}: i grafici servono a semplificare i concetti e a far emergere
immediatamente il dato che interessa, o il dato su cui si basa una parte del
design. Se ne parlerà in lezioni future di documentazione e pitching, ma
intanto un ripasso: esiste una miriade di grafici (``se dico grafico a corde,
qualcuno sa che cos'è?''; il diagramma \emph{sunburst}\dots), alcuni più
esteticamente appaganti di altri. Si può dare colore alle categorie e
modificare qualsiasi cosa, ma bisogna trovare \textbf{il grafico giusto che
comunichi l'idea nel modo ideale}. Per dire la stessa cosa si può usare un
grafico a torta o un grafico a barre; la torta la si può ``bucare'' (ciambella)
o renderla 3D; si può usare un diagramma ad albero (\emph{treemap}), che è
``una torta con altre torte dentro''. Se invece si deve mostrare una
tendenza, ad esempio la progressione delle abilità di un personaggio, serve
una curva. Dipende da che cosa si vuole comunicare: ancora una volta,
\emph{the right tool for the right job}.

\subsection{Colori per separare, aggregare, far emergere}

I \textbf{colori} servono a separare, aggregare e far ``poppare fuori'' le
informazioni (nelle slide del docente, ad esempio, le cose importanti sono in
grassetto e in arancione). Lo stesso vale per tutti gli strumenti di
formattazione: grassetto, sottolineato, corsivo e così via.
\begin{itemize}
  \item \textbf{Coerenza}: se gli attacchi, o un nome, sono scritti in rosso,
        si usa sempre quel rosso.
  \item \textbf{Colori morbidi}, che non distruggano la capacità di vedere le
        cose: un'immagine o una pagina in rosso RGB puro fa ``non vederci più
        per una settimana''.
  \item \textbf{Stampabilità}: un documento con testo bianco su fondo nero ha
        un'ottima visibilità, ma non si può stampare.
\end{itemize}
L'obiettivo è un risultato visivamente appagante, o il più possibile tale.
Come regola del pollice: \textbf{guardare il documento da lontano}, a diverse
percentuali di zoom. Al 50\%, ad esempio, si riesce a identificare il
paragrafo più importante? C'è un testo troppo fitto rispetto alle immagini, o
un'immagine troppo grande rispetto al testo? Bisogna stabilire delle
priorità.

\subsection{Conosci il lettore}

Un documento scritto per un \emph{creative director} è diverso da uno scritto
per un collega. Con il creative director non si sarà estremamente tecnici,
perché non gli interessa, e troppe informazioni creerebbero più problemi che
altro. Per un programmatore ci si preoccupa anche delle \textbf{variabili},
che non interesserebbero a un artista. Se si vuole fare il \emph{tuning}
dell'attacco o della velocità in un sistema di combattimento, all'artista non
interessa; all'animatore forse sì, ma a lui non interessa quanta vita ha il
personaggio o di che colore è la barra. Gli studenti, ad esempio, sanno che il
task che stanno scrivendo va al docente: possono già ragionare secondo alcune
regole, su come ragiona lui, su ciò che vorrebbe vedere un designer che è,
potenzialmente, il loro capo.

\subsection{Sempre aggiornati}

I documenti devono essere \textbf{sempre aggiornati} (\emph{up to date}), un
po' per il processo iterativo, un po' perché la documentazione non resta mai
ferma. Se è ferma, o è il ``sacro graal'' del progetto (la visione), oppure è
\emph{outdated} e qualcosa la sostituirà.

\paragraph{Un grosso asterisco.} La regola è creare la migliore
documentazione possibile affinché il progetto funzioni e venga implementato.
Ma c'è un trucco: \textbf{la migliore documentazione possibile è il progetto
funzionante}. Se si è uno sviluppatore solitario e si ha tutto in testa,
scrivere pagine e pagine o documenti diversi per diverse meccaniche è
probabilmente deleterio: si sa già che cosa fare (ammesso di non cambiare idea
ogni due secondi e di ricordarsi i dati). La documentazione diventa
importante quando bisogna parlare con tante persone, o quando bisogna
\textbf{documentare i passi fatti in passato} per non ripercorrerli. Se si
itera, l'iterazione precedente va documentata: se dovesse tornare un feedback
simile, o qualcuno chiedesse ``perché la meccanica è così e non così?'', si
può rispondere ``perché così l'abbiamo già fatta un mese fa, e questi sono
stati i risultati''. È una tutela per il designer e per il progetto.

Sotto deadline può capitare di non avere il tempo di documentare, o di
documentare bene. In quel caso può convenire non farlo, o farlo ragionando sui
limiti e sulle risorse del progetto: si va direttamente \emph{in game}, e
quando ci sarà più respiro si aggiornerà la documentazione. Almeno il progetto
va avanti. Non è così che si dovrebbe lavorare, ma è il discorso del target
audience applicato ai documenti: chi li legge, e quanto tempo si ha. Se la
deadline è domani, si può fare poco e bisogna tagliare da qualche parte. (Vale,
ad esempio, per le \emph{game jam}, di cui si parlerà.)

\subsection{Iterate}

L'ultimo punto è ormai un mantra: \textbf{iterate}, fin quando si hanno tempo
e risorse.

\begin{table}[htbp]
  \centering
  \small
  \begin{tabular}{@{}p{0.28\textwidth}p{0.64\textwidth}@{}}
    \toprule
    \textbf{Regola} & \textbf{In sintesi} \\
    \midrule
    Mostra, non scrivere & immagini, video, flowchart, link, wireframe
                           (brutti: si itera più in fretta e non si pestano i
                           piedi ad altri dipartimenti) \\
    Essenziale & obiettivo, chiarezza, concisione, niente parole deboli; si
                 procede per sottrazione \\
    Right tool for the right job & slide per presentare, fogli di calcolo per
                                   bilanciare, documenti per analisi e GDD \\
    Coerenza e sequenza & ordine logico, nessuna informazione mancante,
                          grammatica corretta \\
    Mind the language & nomenclatura tecnica, niente parole bandite \\
    Niente decorazioni & nessun elemento inutile, nessuna pagina vuota, non è
                         una tesi \\
    Grafici con intento & il grafico giusto per far emergere il dato giusto \\
    Colori & separare, aggregare, far emergere; coerenti, morbidi,
             stampabili; guardare il documento da lontano \\
    Conosci il lettore & livello di dettaglio diverso per creative director,
                         programmatori, artisti \\
    Sempre aggiornati & tranne la visione; la miglior documentazione è il
                        progetto funzionante \\
    Iterate & finché ci sono tempo e risorse \\
    \bottomrule
  \end{tabular}
  \caption{Le golden rules della documentazione.}
  \label{tab:golden-rules}
\end{table}

% ---------------------------------------------------------------------
\section{Scegliere il documento giusto}
\label{sec:scegliere-documento}

Per ora, come detto, nel Drive si usano \textbf{solo gli strumenti di
Google}: un documento Word inserito nel Drive non viene letto. Le regole per
scegliere il documento sono semplici:

\begin{table}[htbp]
  \centering
  \begin{tabular}{@{}p{0.22\textwidth}p{0.34\textwidth}p{0.34\textwidth}@{}}
    \toprule
    \textbf{Strumento} & \textbf{Quando usarlo} & \textbf{Esempi} \\
    \midrule
    Google Docs & documentazione testuale (\emph{text-telling}) & analisi,
      \emph{game design document} \\
    \addlinespace
    Google Sheets (Excel) & tutto ciò che è numericamente pesante & matrici da
      bilanciare, grafici da costruire e poi importare nei documenti; è ``casa''
      degli \emph{economy designer} \\
    \addlinespace
    Google Slides (PowerPoint) & comunicazione estremamente visuale & qualsiasi
      presentazione, qualsiasi \emph{pitch} \\
    \bottomrule
  \end{tabular}
  \caption{Quale strumento della suite di Google usare.}
  \label{tab:strumenti-google}
\end{table}

Esistono alternative, e dipendono da dove si lavora. Il docente, ad esempio,
in Stormind Games ha usato molti metodi diversi a seconda dei progetti:
uno richiedeva Confluence, un altro Miro, e ogni volta bisognava adattarsi a
quel supporto. Nel corso ci si adatta agli strumenti indicati dal docente,
perché \textbf{le aziende fanno così}: danno dei limiti, e si lavora entro
quei limiti.

\subsection{Less is more}

Meno si scrive, con più cognizione di causa, più si taglia, e più si diventa
comprensibili.

\paragraph{Il documento di Diablo.} Il docente mostra una pagina di un vecchio
documento di design e chiede di indovinare il gioco senza leggere: un
personaggio ``demoniaco, maligno''. È \emph{Diablo}: uno dei documenti di
design del primo \emph{Diablo}, del 1994.\footnote{Si tratta del documento di
\emph{pitch} di \emph{Diablo} redatto da Condor (poi Blizzard North) nel
1994, oggi liberamente consultabile online.} Da allora c'è stata
un'evoluzione nel modo di fare documentazione, grazie ad alcuni ``padri'' che
hanno cambiato il modo di documentare, tra cui \textbf{Stone Librande}: si è
arrivati al punto in cui la documentazione può ridursi a \textbf{un singolo
documento di una pagina} (non onnicomprensivo, ma il più possibile tale; su
questo concetto si tornerà).\footnote{Stone Librande ha presentato il
proprio approccio dei \emph{one-page designs} alla Game Developers Conference
del 2010.}

\paragraph{Le design bible.} In passato (``sfortunatamente passato'') si
facevano le cosiddette \textbf{bibbie di design}: enormi tomi da 50, 100, 150
pagine che condensavano tutto il gioco e che, teoricamente, tutti dovevano
leggere, ma che nessuno leggeva davvero (``ma non te lo leggerà mai!''). Si
facevano per ignoranza, che ci sta, e perché il mondo del design videoludico
va molto veloce e non sempre tiene il passo con certe modifiche. Ancora oggi,
ad esempio, alcune di queste idee vengono richieste per ottenere fondi: uno
dei fondi europei richiede che il gioco abbia la \emph{design bible}, un
documento enorme di tutto il gioco. Ma nessuno parte da lì: la si fa perché si
vogliono quei soldi, per ``smarcare'' l'obbligo; la compone il producer,
incollando insieme i documenti che gli vengono dati. Si è andati molto oltre:
così come oggi nessuno leggerebbe una bibbia di design (e in realtà nessuno
la leggeva nemmeno allora; ``non so quanti alberi siano stati tagliati per
bibbie di design da mettere su ventimila scrivanie''), l'obiettivo è
raggiungere qualcosa di \textbf{efficace e condensato} in molte meno pagine,
e diviso anche per dipartimento (se ne parlerà più avanti).

\paragraph{La documentazione è ovunque.} La comunicazione e la
documentazione arrivano dappertutto. La documentazione di \textbf{Unreal} è
documentazione tecnica a tutti gli effetti: documenti localizzati, una
\emph{wiki}. Quando serve qualcosa, ad esempio informazioni sulle
\emph{nav mesh}, la si cerca lì, e quella documentazione permette poi di
lavorare sui blueprint: è un ponte che aiuta a fare cose, a implementare. Va
capito quanto sia efficace e quanti problemi risolva, ma il concetto è che
\textbf{più è efficace, meglio è}.

\subsection{Strumenti aggiuntivi}

Alcuni strumenti utili, soprattutto per rispondere alla regola \emph{mostra,
non scrivere} (non obbligatori, ma buoni per un lavoro preciso):
\begin{description}
  \item[Photopea] alternativa gratuita a Photoshop (che costa), utilizzabile
        online, per modificare foto e immagini e fare \emph{mash-up}; ci sono
        le pubblicità, ma sono ignorabili.
  \item[Lucidchart] (a pagamento) e \textbf{draw.io} (gratuito) per creare
        \emph{chart}, flowchart, wireframe e qualsiasi altra cosa: i designer
        ci fanno anche dei template. ``Uno strumentopolo misterioso che ci
        servirà più avanti''.
  \item[Miro] gratuito con dei limiti, per raccogliere idee ma anche,
        potenzialmente, per documentare.
  \item[Paint] il ``buon caro amico'' Paint. Se si sta discutendo con un
        designer e bisogna spiegargli qualcosa in maniera immediata e chiara,
        uno schizzo su Paint ``e passa la paura''. Uno studente obietta che si
        può disegnare con qualsiasi strumento: certo, ma Paint è
        \emph{percepito} come più grezzo, e quindi, in chiamata, si capisce
        che è una cosa destinata al cestino; si perde anche quell'attaccamento
        all'idea.
\end{description}

\begin{quote}
\centering
\textbf{Meglio brutto e chiaro che perfetto, preciso, con tutte le
spaziature, ma non chiaro e potenzialmente fraintendibile.}
\end{quote}
(Il docente precisa che lo dice anche per sé.)

% ---------------------------------------------------------------------
\section{Le regole formali del corso}
\label{sec:regole-formali}

In chiusura il docente aggiunge ``un paio di trucchi, un paio di catene e
catenacci'': le regole formali che valgono per tutti i documenti del corso
(verranno aggiunte anche al documento delle regole; ``non fate i furbi'').

\begin{description}
  \item[Margini] tutti i documenti di Google Docs devono avere i margini
        \textbf{1,5 sopra, 1,5 sotto, 1,25 a sinistra, 1,25 a
        destra}.\footnote{L'unità di misura non viene specificata in lezione;
        nell'interfaccia italiana di Google Docs i margini sono espressi in
        centimetri.} Il documento deve essere \textbf{stampabile}: niente
        documenti con lo sfondo nero.
  \item[Niente formato senza pagine] Google Docs permette di usare un formato
        ``senza pagine'', a pagina singola continua: è \textbf{bandito}, per
        ora. Non aiuta né nella stampabilità né nel ragionare entro i limiti
        della documentazione.
  \item[Font] il font è categoricamente \textbf{Roboto 11}, per tutti i testi.
        Titoli e intestazioni sono esclusi dalla regola: anche lì il docente
        consiglia caldamente Roboto, ma si è più liberi. Se nel testo c'è
        anche una sola parola che non è in Roboto 11, il documento non viene
        letto (``state sicuri che io la becco''). Uno studente chiede se si
        possa impostare Roboto come predefinito: sì, si può fare.
  \item[Sfruttare la toolbar] grassetto, corsivo, sottolineato, interlinea,
        interruzioni: gli strumenti sono moltissimi.
  \item[Colori banditi] i colori rientrano negli strumenti, ma alcuni sono
        banditi. Il docente non ne fa la lista (``siete abbastanza adulti per
        riconoscerli''): è chiaro che i colori RGB puri non funzionano, che il
        blu acceso si legge a fatica, che il giallo su bianco è illeggibile,
        che il testo nero su rosso ``fa perdere due o tre diottrie''.
  \item[Link] il sottolineato non è necessariamente un link, ma un link deve
        essere \textbf{ben differenziato} da tutto il resto: se non si sa che
        lì c'è un link, non ci si clicca mai; viceversa, se si pensa che ogni
        sottolineato sia un link, si clicca ovunque senza andare da nessuna
        parte.
  \item[Rinominare tutto] il documento, il titolo e qualsiasi cosa contenga
        che sia navigabile, come le \emph{tab}. Vale per Google Docs, per
        Google Sheets (ogni foglio va rinominato) e per qualsiasi altro
        strumento.
  \item[Margini e immagini] un'immagine che sfora i margini, un'immagine
        sopra un'altra immagine, un testo che si sovrappone a un'immagine o
        un'immagine che copre parte del testo: il documento non viene letto.
  \item[Coerenza] evitare incoerenze ed errori grammaticali. Se si chiude con
        il punto ogni frase o ogni voce di un elenco puntato, lo si fa per
        tutte. Chiudere le voci con il punto e virgola e l'ultima con il punto
        è accettato (``non sono un grande fan, ma è una questione di
        preferenze''), purché sia coerente; non lo è ``punto, punto e
        virgola, punto e virgola, punto''.
\end{description}

\paragraph{Si può fare tutto.} Il mantra da portarsi a casa:
\textbf{con i documenti di Google si può fare tutto}. Sembrano uno strumento
limitante, ma ``si può fare il mondo''. Non lo si imparerà subito: lo si
imparerà piano piano, soprattutto vedendo qualcuno che ha fatto qualcosa che
non si pensava possibile, o vedendolo fare al docente. L'idea non è copiare,
ma \textbf{capire}, mettere la tecnica nel proprio bagaglio e sapere quando
usarla. Il docente, ad esempio, è un fan delle tabelle solo fino a un certo
punto, mentre i colleghi di Roma ne sono grandissimi fan, ``fin troppo''. Non
è che, se qualcuno è un maestro delle tabelle, allora si usano sempre quelle:
si imparano i punti di forza degli altri, si prendono i propri e si migliora;
gli altri vedono e migliorano a loro volta.

Gli strumenti da conoscere includono le \textbf{tabelle} (con il loro spazio
interno: si possono allungare, colorare, se ne possono togliere le linee), gli
\textbf{elenchi puntati} (difficile trovare una lista più leggibile) e la
gestione degli \textbf{spazi}. Lo spazio è prezioso, e cederlo volutamente è
rischioso, ma può portare grandi risultati: come sempre, a un rischio
corrisponde un \emph{reward}.

% ---------------------------------------------------------------------
\section{Il task della settimana}
\label{sec:task-gd3}

Non c'è un task nuovo: il task resta quello sui \emph{core elements} dei
giochi arcade (sezione~\ref{sec:task-gd2}), con la stessa scadenza
(\textbf{mercoledì alle 13}). Ma ora va svolto applicando le regole viste in
questa lezione.
\begin{enumerate}
  \item Riguardare il documento già iniziato seguendo le \textbf{golden rules}
        (sezione~\ref{sec:golden-rules}).
  \item Applicare le \textbf{regole formali} (sezione~\ref{sec:regole-formali}):
        margini, Roboto 11, colori, coerenza.
  \item Il documento consegnato dovrà rispettare entrambe.
\end{enumerate}

\paragraph{Come procedere: la risposta a Federica.} Federica chiede come
collegare concretamente queste regole al task, rendendo il documento
accattivante con le immagini giuste. La risposta non è semplice: obbliga a
\textbf{pensare}, dopo la lezione, a ciò che si è detto, a rimuginarci, o a
riguardare la lezione. Alcune indicazioni concrete:
\begin{itemize}
  \item i \textbf{margini} vanno applicati a qualsiasi documento si stia
        facendo;
  \item uno studente di Roma ha impostato il task come una serie di slide:
        probabilmente non è il supporto adatto, perché si tratta di analisi e
        ricerca sui core elements, e il miglior candidato è un documento di
        Google. Un altro documento mostrato usa correttamente Roboto 11, ma non
        ha i margini corretti. Conviene quindi trasporre il lavoro già fatto
        in un documento di Google con i margini e il testo corretti, e
        lavorare lì;
  \item prendere le informazioni che si hanno e provare a modificarle:
        \emph{come} sono scritte, \emph{quante} sono, che cosa si può
        \emph{mostrare}. Ad esempio, se si elencano i \emph{game piece} di
        Schell, si può usare un elenco puntato; e se esistono gli sprite di
        quei pezzi di gioco, può essere utile mostrarli;
  \item usare lo \textbf{stesso colore} per parlare della stessa cosa, ad
        esempio dello stesso pezzo di gioco;
  \item usare il \textbf{linguaggio tecnico} ed evitare le parole bandite;
  \item imparare a non sprecare lo spazio (``lo vedrete sbagliando''); e per i
        colori, se si ha il dubbio che un colore non sia leggibile (come un
        verde chiaro su bianco), probabilmente non lo è.
\end{itemize}
È difficile assorbire tutto e subito: ci vogliono tempo, esperienza e tanta
pratica. ``Nessun documento andrà bene'', rassicura il docente: è il primo
documento che si consegna, e sarebbe strano il contrario. ``Imparerò
sbagliando'', conclude Federica: ``così si fa''.

\paragraph{Lasciar sedimentare.} Un ultimo consiglio, non scritto nelle
slide: \textbf{non fare tutto subito}. Fatto il documento, lo si lascia
riposare per qualche ora o un giorno, senza toccarlo, e lo si riguarda il
giorno dopo: sicuramente alcune cose appariranno in modo diverso. Anche questo
fa parte del processo iterativo. I designer non bilanciano a ruota tutto il
giorno fino ad arrivare al punto corretto: se fanno lo stesso lavoro per tutto
il giorno, si \textbf{abituano all'errore}, a qualcosa che magari è sbagliato
ma che ormai sono abituati a vedere (un salto di quell'altezza, ad esempio).
Di solito si lavora un po' oggi, ci si ferma, si riprende domani e si verifica
se il \emph{feeling} è lo stesso del giorno prima. La settimana a disposizione
va sfruttata così: non lavorando ogni ora di ogni giorno (irrealistico), ma
dedicando quotidianamente un po' di tempo al task.

\paragraph{Prossimi appuntamenti.}
\begin{itemize}
  \item La lezione di lunedì si terrà in presenza; ricordarsi di chiedere al
        docente il concetto di Raph Koster rimasto in sospeso.
  \item Dopo la consegna e i commenti, entro la settimana successiva
        (al massimo venerdì), si farà una \textbf{dimostrazione}: si prenderà
        un documento e lo si formatterà da capo, per vedere qualche trucco.
        Poiché venerdì è previsto uno sciopero dei mezzi, quella lezione
        potrebbe tenersi da remoto.
  \item Per dubbi e domande si possono usare le chat o parlarne a voce; ma a
        domande come ``questo colore è corretto?'' la risposta sarà ``ci
        rivediamo mercoledì''.
\end{itemize}

% ---------------------------------------------------------------------
\section*{Punti chiave del capitolo}
\addcontentsline{toc}{section}{Punti chiave del capitolo}

\begin{itemize}
  \item Nei framework di core elements: i \textbf{game piece} sono oggetti che
        cambiano stato, i \textbf{verbi} sono ciò che il giocatore può fare,
        le \textbf{regole} sono tutto il resto. I \textbf{player} di Romero e
        Schreiber sono il target audience e vanno definiti con dati.
  \item Lo stesso termine può avere significati diversi in framework diversi
        (\emph{estetiche}, \emph{mastery}\dots): l'obiettivo è parlare la
        stessa lingua del team.
  \item Con Bura si individua la cella di una meccanica e la si potenzia o
        blocca per ottenere un responso emotivo; il designer percorre il ciclo
        variabili--emozioni--interazioni \textbf{a ritroso}.
  \item Il ciclo di produzione comprende preproduzione, produzione e rilascio;
        per i videogiochi si parla di \textbf{DevOps} (prima del rilascio) e
        \textbf{LiveOps} (dopo). Più è profonda la preproduzione, meglio
        vanno le altre fasi.
  \item La parte più importante del ciclo è \textbf{la freccia che torna
        indietro}: tutti i cicli di produzione sono \textbf{iterativi}, e
        fare design iterativo è l'unico modo per avere un buon prodotto.
  \item Ogni decisione di design è fatta per soddisfare il \textbf{target
        audience}; il gioco è di chi lo riceve, non di chi ci lavora.
  \item L'iterazione è \textbf{involvement}, \textbf{iteration},
        \textbf{holistic} e \textbf{adaptive}; i \textbf{pilastri} invece non
        si cambiano.
  \item I due segreti: \textbf{più lungo è il processo iterativo, migliore è il
        risultato}; \textbf{prima createlo, poi fatelo diventare bello}.
  \item La documentazione è il primo lavoro del designer e la sua firma:
        mantiene coerente la visione, aggiorna, supporta, controlla, ed
        \textbf{evolve} con il progetto. Il designer è Mercy, non Genji.
  \item Golden rules: mostra non scrivere, essenziale, strumento giusto,
        coerenza, linguaggio tecnico, niente decorazioni, grafici con intento,
        colori con criterio, conosci il lettore, sempre aggiornati, itera.
  \item Regole del corso: solo suite di Google, margini 1,5/1,5/1,25/1,25,
        Roboto 11, documenti stampabili, tutto rinominato, coerenza formale.
\end{itemize}
